
\section{Group homology/cohomology}

\subsection{Types of groups}

\begin{definition} \label{definition:discrete_group}
A \hldef{discrete group} is a \CrefAndHyperrefIfExist{definition:group}{group} $G$ equipped with the \CrefAndHyperrefIfExist{definition:discrete_topological_space_of_a_set}{discrete topology}.
Under this \CrefAndHyperrefIfExist{definition:topological_space}{topology}, every singleton set $\{g\}$ for $g \in G$ is \CrefAndHyperrefIfExist{definition:topological_space}{open}, and the group operations (multiplication $G \times G \to G$ and inversion $G \to G$) are necessarily \CrefAndHyperrefIfExist{definition:continuous_map_between_open_subsets_of_euclidean_spaces}{continuous}.
\end{definition}

\begin{proposition} \label{proposition:topological_group_is_discrete_if_and_only_if_identity_element_forms_an_open_set}
Let $G$ be a \CrefAndHyperrefIfExist{definition:topological_group}{topological group}. 
$G$ is a \CrefAndHyperrefIfExist{definition:discrete_group}{discrete group} if and only if the singleton set $\{e\}$, where $e \in G$ is the \CrefAndHyperrefIfExist{definition:left_right_identity_element_of_a_binary_operation_on_a_set}{identity element}, is an \CrefAndHyperrefIfExist{definition:topological_space}{open} set.
\end{proposition}


\subsection{Group ring of a group over a ring}


\begin{definition} \label{definition:group_ring_of_a_group_over_a_ring}
Let $G$ be a \CrefAndHyperrefIfExist{definition:group}{group} and $R$ a \CrefAndHyperrefIfExist{definition:ring}{not necessarily commutative ring}.

The \hldef{group ring} \hl{$R[G]$}, also denoted \hl{$RG$}, is the set of all formal linear combinations of the form
$$ \sum_{g \in G} r_g g $$
where $r_g \in R$ and $r_g = 0$ for all but finitely many $g \in G$. This set is endowed with a ring structure via:
\begin{itemize}
    \item Addition:
    $$ \left( \sum_{g \in G} r_g g \right) + \left( \sum_{g \in G} s_g g \right) = \sum_{g \in G} (r_g + s_g) g $$
    \item Multiplication: Defined by the extension of the group law and the requirement that $R$ and $G$ commute ($rg = gr$):
    $$ \left( \sum_{g \in G} r_g g \right) \left( \sum_{h \in G} s_h h \right) = \sum_{k \in G} \left( \sum_{gh=k} r_g s_h \right) k $$
\end{itemize}
The identity element of $R[G]$ is $1_R 1_G$.

$R[G]$ may also be called the \hldef{group algebra of $G$ over $R$} as the natural ring homomorphism $R \to R[G]$ sending $r \mapsto r 1_G$ is an \CrefAndHyperrefIfExist{definition:algebra_of_a_ring}{algebra}
\end{definition}
\begin{definition} \label{definition:augmentation_map_of_a_group_ring_to_its_base_ring}
Let $G$ be a \CrefAndHyperrefIfExist{definition:group}{group} and $R$ a \CrefAndHyperrefIfExist{definition:ring}{not necessarily commutative ring}.
The \hldef{augmentation map} \hl{$\epsilon: R[G] \to R$} is the $R$-linear \CrefAndHyperrefIfExist{definition:ring_homomorphism}{ring homomorphism} defined by $\epsilon(g) = 1$ for all $g \in G$. The kernel of this map, denoted by notations such as \hl{$I_G$}, is called the \hldef{augmentation ideal}.
\end{definition}
\begin{definition} \label{definition:norm_element_of_a_group_algebra_of_a_finite_group_over_a_ring}
    Let $R$ be a \CrefAndHyperrefIfExist{definition:ring}{ring}. Let $G$ be a \CrefAndHyperrefIfExist{definition:countable_finite_uncountable_sets}{finite} \CrefAndHyperrefIfExist{definition:group}{group}. 
    The \hldef{norm elements} of the \CrefAndHyperrefIfExist{definition:group_ring_of_a_group_over_a_ring}{group ring} $R[G]$ is the sum $\sum_{g \in G} g$. 
\end{definition}
\begin{definition} \label{definition:action_of_a_group_on_an_abelian_group}
Let $G$ be a group and $A$ an abelian group. 

\begin{enumerate}

    \item An \hldef{action of $G$ on $A$ on the left} is a map $\cdot: G \times A \to A$ such that for all $g, h \in G$ and $a, b \in A$:

    \begin{itemize}
        \item $g \cdot (a + b) = (g \cdot a) + (g \cdot b)$ (the action is by automorphisms of $A$)
        \item $(gh) \cdot a = g \cdot (h \cdot a)$
        \item $1_G \cdot a = a$
    \end{itemize}

    \item An \hldef{action of $G$ on $A$ on the right} is a map $\cdot: A \times G \to A$ such that for all $g, h \in G$ and $a, b \in A$:
    \begin{itemize}
        \item $(a + b) \cdot g = (a \cdot g) + (b \cdot g)$ (the action is by automorphisms of $A$)
        \item $a \cdot (gh) = (a \cdot g) \cdot h$
        \item $a \cdot 1_G = a$
    \end{itemize}
\end{enumerate}
In particular, an action of $G$ on $A$ is an \CrefAndHyperrefIfExist{definition:left_and_right_group_actions_on_a_set}{action of $G$} on the underlying set of $A$ with extra conditions.

When such an action is specified, $A$ is called a \hldef{$G$-module} (or more specifically, a \hldef{$\mathbb{Z}[G]$-module}). 

\TextIfExists{definition:left_right_module_over_a_group_ring_of_a_group_over_a_ring}{Note that this definition is a special case of \Cref{definition:left_right_module_over_a_group_ring_of_a_group_over_a_ring} in the case that the base ring $R$ is $\bbZ$.}
\end{definition}
\begin{definition} \label{definition:left_right_module_over_a_group_ring_of_a_group_over_a_ring}
Let $G$ be a \CrefAndHyperrefIfExist{definition:group}{group} and $R$ a \CrefAndHyperrefIfExist{definition:ring}{not necessarily commutative ring}. 
Write $R[G]$ for the associated \CrefAndHyperrefIfExist{definition:group_ring_of_a_group_over_a_ring}{group ring}.
\begin{enumerate}
    \item A \hldef{left $R[G]$-module} or a \hldef{$G$-module over $R$ on the left} is a \CrefAndHyperrefIfExist{definition:module_of_a_ring}{left $R$-module} $A$ equipped with a left \CrefAndHyperrefIfExist{definition:action_of_a_group_on_an_abelian_group}{$G$-action} such that for all $g \in G$, $r \in R$, and $a, a' \in A$:
    \begin{itemize}
        \item $g(ra) = r(ga)$ (The $G$-action is $R$-linear).
        \item $g(a + a') = ga + ga'$ (The $G$-action is additive).
    \end{itemize}

    \item A \hldef{right $R[G]$-module} or a \hldef{$G$-module over $R$ on the right} is a right \CrefAndHyperrefIfExist{definition:module_of_a_ring}{$R$-module} $A$ equipped with a right $G$-action such that for all $g \in G$, $r \in R$, and $a, a' \in A$:
    \begin{itemize}
        \item $(ar)g = (ag)r$ (The $G$-action is $R$-linear).
        \item $(a + a')g = ag + a'g$ (The $G$-action is additive).
    \end{itemize}
\end{enumerate}
Equivalently, a left (resp. right) $R[G]$-module is an \CrefAndHyperrefIfExist{definition:group}{abelian group} $A$ together with a ring homomorphism $\rho: R[G] \to \text{End}_{\mathbb{Z}}(A)$ (resp. $\rho: R[G]^{op} \to \text{End}_{\mathbb{Z}}(A)$).

Equivalently, a left/right $R[G]$-module is a \CrefAndHyperrefIfExist{definition:module_of_a_ring}{left/right module} of the group ring $R[G]$.

A \hldef{left/right $G$-module} usually refers to a left/right $\bbZ[G]$-module.

A $G$-module over $R$ is equivalent to a \CrefAndHyperrefIfExist{definition:representation_of_a_group_on_a_module_over_a_not_necessarily_commutative_ring}{representation of $G$ on an $R$-module}. 
\end{definition}

\subsection{Invariant/coinvariant functors and Group homology/cohomology}
\begin{definition} \label{definition:invariant_subgroup_of_a_group_acting_on_an_abelian_group}

Let $A$ be an \CrefAndHyperrefIfExist{definition:group}{abelian group} with a \CrefAndHyperrefIfExist{definition:action_of_a_group_on_an_abelian_group}{$G$-action}, say on the left. 
The \hldef{group of $G$-invariant elements}, denoted \hl{$A^G$} (this notation is also standard when $G$ \CrefAndHyperrefIfExist{definition:action_of_a_group_scheme_on_a_scheme}{acts on} $A$ on the right), is \CrefAndHyperrefIfExist{definition:subgroup_of_a_group}{subgroup} of $A$ defined as:
$$ A^G = \{ a \in A \mid g \cdot a = a \text{ for all } g \in G \}.$$
\TextIfExists{definition:fixed_point_set_of_a_group_acting_on_a_set}{Note that this coincides with the fixed point sete of $G$ in $A$ as per \Cref{definition:fixed_point_set_of_a_group_acting_on_a_set}.}

If $A$ is an \CrefAndHyperrefIfExist{definition:left_right_module_over_a_group_ring_of_a_group_over_a_ring}{(left/right) $R[G]$-module} where $R[G]$ is the \CrefAndHyperrefIfExist{definition:group_ring_of_a_group_over_a_ring}{group ring}, then $A^G$ is an (left/right) \CrefAndHyperrefIfExist{definition:module_of_a_ring}{$R$-module}. In fact, $A^G$ is naturally regarded as a \CrefAndHyperrefIfExist{definition:trivial_G_module_of_a_group_over_a_ring}{trivial $R[G]$-module}.
\end{definition}
\begin{definition} \label{definition:coinvariant_quotient_group_of_a_group_acting_on_an_abelian_group}
Let $A$ be an \CrefAndHyperrefIfExist{definition:group}{abelian group} with a \CrefAndHyperrefIfExist{definition:action_of_a_group_on_an_abelian_group}{$G$-action}, say on the left. 
The of \hldef{group of $G$-coinvariant elements}, denoted \hl{$A_G$} (this notation is also standard when $G$ \CrefAndHyperrefIfExist{definition:action_of_a_group_scheme_on_a_scheme}{acts on} $A$ on the right), is the \CrefAndHyperrefIfExist{definition:quotient_of_a_group_by_a_normal_subgroup}{quotient group} defined as:
$$ A_G = A / \langle \{ (ga-a): a \in A, g \in G \} \rangle .$$
\CrefIfExists{definition:ideal_generated_by_a_subset_of_a_ring}
If $A$ is an \CrefAndHyperrefIfExist{definition:left_right_module_over_a_group_ring_of_a_group_over_a_ring}{(left/right) $R[G]$-module} where $R[G]$ is the \CrefAndHyperrefIfExist{definition:group_ring_of_a_group_over_a_ring}{group ring}, then $A_G$ is an (left/right) \CrefAndHyperrefIfExist{definition:bimodule_over_rings_module_over_a_ring}{$R$-module}.
In fact, $A_G$ is naturally regarded as a \CrefAndHyperrefIfExist{definition:trivial_G_module_of_a_group_over_a_ring}{trivial $R[G]$-module}.
\end{definition}
\begin{definition} \label{definition:trivial_G_module_of_a_group_over_a_ring}
Let $G$ be a \CrefAndHyperrefIfExist{definition:group}{group} and $R$ a \CrefAndHyperrefIfExist{definition:ring}{not necessarily commutative ring}.
\begin{enumerate}
    \item A \hldef{trivial $G$-module over $R$} or a \hldef{trivial $R[G]$-module} is an \CrefAndHyperrefIfExist{definition:left_right_module_over_a_group_ring_of_a_group_over_a_ring}{$R[G]$-module} (say left) $A$ on which $G$ acts trivial, i.e. $ga = a$ for all $g \in G$ and $a \in A$.

    \item Given a (left/right) $R$-module $A$, its \hldef{trivial $G$-module (over $R$)} is the $R[G]$-module structure whose underlying $R$-module is $A$ and whose $G$-action is trivial. This construction yields \CrefAndHyperrefIfExist{definition:additive_functor_between_additive_categories}{additive functors}
    $$_{R} \Mod \to _{R[G]} \Mod$$
    $$\Mod_{R} \to \Mod_{R[G]}$$
    \CrefIfExists{definition:category_of_modules_and_bimodules_over_rings}
\end{enumerate}
\end{definition}
\begin{lemma}[cf. {\cite[Exercise 6.1.1]{weibel}}] \label{lemma:invariant_coinvariant_functors_for_group_rings_are_adjoint_to_trivial_module_functor}
    Let $G$ be a \CrefAndHyperrefIfExist{definition:group}{group} and $R$ a \CrefAndHyperrefIfExist{definition:ring}{not necessarily commutative ring}. Write $R[G]$ for the \CrefAndHyperrefIfExist{definition:group_ring_of_a_group_over_a_ring}{group ring} and let $A$ be a \CrefAndHyperrefIfExist{definition:left_right_module_over_a_group_ring_of_a_group_over_a_ring}{(left) $R[G]$-module}. 

    \begin{enumerate}
        \item \CrefAndHyperrefIfExist{definition:invariant_subgroup_of_a_group_acting_on_an_abelian_group}{$A^G$} is the maximal \CrefAndHyperrefIfExist{definition:trivial_G_module_of_a_group_over_a_ring}{trivial} \CrefAndHyperrefIfExist{definition:submodule_of_a_module_over_a_ring}{submodule} of $A$. Moreover, $-^G: {}_{R[G]}\Mod \to {}_{R}\Mod$\CrefIfExists{definition:category_of_modules_and_bimodules_over_rings} is \CrefAndHyperrefIfExist{definition:adjoint_functors_between_categories_unit_counit_of_adjoint_functors}{right adjoint} to the \CrefAndHyperrefIfExist{definition:trivial_G_module_of_a_group_over_a_ring}{trivial module functor} $_{R}\Mod \to {}_{R[G]}\Mod$. 

        In particular, $-^G$ is \CrefAndHyperrefIfExist{definition:exact_functor_between_abelian_categories}{left exact} (\Cref{proposition:left_right_adjoint_of_add_cats_is_add_and_right_left_exact_if_between_ab_cats}).

        \item \CrefAndHyperrefIfExist{definition:coinvariant_quotient_group_of_a_group_acting_on_an_abelian_group}{$A_G$} is the maximal \CrefAndHyperrefIfExist{definition:trivial_G_module_of_a_group_over_a_ring}{trivial} \CrefAndHyperrefIfExist{definition:quotient_module_of_a_module_by_a_module_of_a_ring}{quotient module} of $A$. Moreover, $-_G: {}_{R[G]}\Mod \to {}_{R}\Mod$\CrefIfExists{definition:category_of_modules_and_bimodules_over_rings} is \CrefAndHyperrefIfExist{definition:adjoint_functors_between_categories_unit_counit_of_adjoint_functors}{left adjoint} to the \CrefAndHyperrefIfExist{definition:trivial_G_module_of_a_group_over_a_ring}{trivial module functor} $_{R}\Mod \to {}_{R[G]}\Mod$.

        In particular, $-_G$ is \CrefAndHyperrefIfExist{definition:exact_functor_between_abelian_categories}{right exact} (\Cref{proposition:left_right_adjoint_of_add_cats_is_add_and_right_left_exact_if_between_ab_cats}).
    \end{enumerate}
\end{lemma}
\begin{proof}
    \begin{enumerate}
        \item By definition, $A^G$ consists of all $a \in A$ such that $ga = a$. Therefore, $A^G$ is the maximal trivial $R[G]$-submodule of $A$.

        Let $M$ be any $R$-module and let $N$ be any $R[G]$-module. By \Cref{lemma:invariant_coinvariant_functors_for_group_rings_are_naturally_isomorphic_to_tensor_hom_by_trivial_module}, there is a \CrefAndHyperrefIfExist{definition:natural_transformation_between_functors_between_categories}{natural isomorphism} $N^G \cong \Hom_{R[G]}(R,N)$. Regarding $M$ as a trivial $R[G]$-module, we have a natural isomorphism
        $$\Hom_{R[G]}(M, N) \cong \Hom_{R[G]}(R \otimes_R M, N) \cong \Hom_R(M, \Hom_{R[G]}(R,N))$$
        where $R$ is given the structure of a trivial left $R[G]$-module. By \CrefAndHyperrefIfExist{theorem:tensor_hom_adjunction_for_bimodules_of_rings}{tensor-hom adjunction}, this is naturally isomorphic to 
        $$\Hom_{R}(M, N^G)$$

        \item Clearly, $A_G$ is a trivial $R[G]$-module. Moreover, for any quotient module of $A$ on which $G$ acts trivially, we must have that $ga$ and $a$ are identified in that quotient module for any $a \in A$ and $g \in G$. Thus, any quotient of $A$ that is a trivial $G$-module must be a quotient of $A_G$.

        Let $M$ be any (left) $R[G]$-module and let $N$ be any (left) $R$-module. By \Cref{lemma:invariant_coinvariant_functors_for_group_rings_are_naturally_isomorphic_to_tensor_hom_by_trivial_module}, there is a \CrefAndHyperrefIfExist{definition:natural_transformation_between_functors_between_categories}{natural isomorphism} $M_G \cong M \otimes_{R[G]} R$ where $R$ is equipped with a $R$-$R[G]$-bimodule structure with the $R[G]$-module structure trivial. Therefore, we have a natural isomorphism
        $$\Hom_R(M_G, N) \cong \Hom_{R[G]}(M, N)$$
        where $N$ has the structure of a trivial (left) $R[G]$-module. The left hand side of the above natural isomorphism is thus identifiable as 
        $$\Hom_R(M \otimes_{R[G]} R, N)$$
        which is naturally isomorphic to 
        $$\Hom_{R[G]}(M, \Hom_{R}(R, N))$$
        by the tensor-hom adjunction. Since $G$ acts trivially on $R$, the induced $R[G]$-module structure on $\Hom_R(R, N)$ is also trivial: for any $\phi \in \Hom_R(R, N)$, $(g \cdot \phi)(r) = \phi(r \cdot g) = \phi(r)$.  Therefore, $R$-module $N$ is naturally isomorphic to $\Hom_R(R, N)$. We thus have a natural isomorphism 
        $$\Hom_{R[G]}(M, \Hom_{R}(R, N)) \cong \Hom_{R[G]}(M, N)$$
        which shows the desired adjunction
        $$\Hom_R(M_G, N) \cong \Hom_{R[G]}(M, N)$$

    \end{enumerate}
\end{proof}
\begin{lemma}[cf. {\cite[Lemma 6.1.1]{weibel}}] \label{lemma:invariant_coinvariant_functors_for_group_rings_are_naturally_isomorphic_to_tensor_hom_by_trivial_module}
    Let $G$ be a \CrefAndHyperrefIfExist{definition:group}{group} and $R$ a \CrefAndHyperrefIfExist{definition:ring}{not necessarily commutative ring}. Write $R[G]$ for the \CrefAndHyperrefIfExist{definition:group_ring_of_a_group_over_a_ring}{group ring} and let $A$ be a \CrefAndHyperrefIfExist{definition:left_right_module_over_a_group_ring_of_a_group_over_a_ring}{left $R[G]$-module}. 

    Consider $R$ as a \CrefAndHyperrefIfExist{definition:trivial_G_module_of_a_group_over_a_ring}{trivial $G$-module over $R$}; one may in fact obtain a \CrefAndHyperrefIfExist{definition:module_of_a_ring}{two-sided $R[G]$ module structure} in doing so. We have \CrefAndHyperrefIfExist{definition:natural_transformation_between_functors_between_categories}{natural isomorphisms}
    $$A_G \cong R \otimes_{R[G]} A$$
    (\Cref{definition:coinvariant_quotient_group_of_a_group_acting_on_an_abelian_group}) \CrefIfExists{definition:tensor_product_of_bimodules_of_rings}
    of left $R$-modules by considering $R$ as an $R$-$R[G]$-bimodule and 
    $$A^G \cong \Hom_{R[G]}(R, A)$$
    (\Cref{definition:invariant_subgroup_of_a_group_acting_on_an_abelian_group}) \CrefIfExists{definition:hom_of_left_right_bi_modules_of_rings}
    of left $R$-modules by considering $R$ as an $R[G]$-$R$-bimodule.
\end{lemma}

\begin{proof}
    \TODO{}
    Consider the functor
    $$\Hom_R(R, -): _{R}\Mod \to _{R[G]} \Mod$$
    where $\Hom_R$ here is the group of left $R$-module homomorphisms and $R$ is considered as an $R-R[G]$-bimodule.  
\end{proof}

\begin{definition} \label{definition:group_cohomology_homology_of_a_module_over_a_group_ring_over_a_ring}
Let $G$ be a \CrefAndHyperrefIfExist{definition:group}{group} and $R$ a \CrefAndHyperrefIfExist{definition:ring}{not necessarily commutative ring}. Let $A$ be a \CrefAndHyperrefIfExist{definition:left_right_module_over_a_group_ring_of_a_group_over_a_ring}{left/right $R[G]$-module}.

\begin{enumerate}
    \item The \hldef{cohomology group of $G$ with coefficients in $A$} or the \hldef{group cohomology of $A$ as an $R[G]$-module}, denoted \hl{$H^n(G, A)$} for $n \geq 0$, is the $n$-th \CrefAndHyperrefIfExist{definition:left_right_derived_functors_of_a_right_left_exact_functor_between_abelian_categories_where_source_has_enough_projectives_injectives}{right derived functor} object $(R^n(-^G))(A)$ of the \CrefAndHyperrefIfExist{definition:invariant_subgroup_of_a_group_acting_on_an_abelian_group}{invariants} functor $A \mapsto A^G$, which is \CrefAndHyperrefIfExist{definition:exact_functor_between_abelian_categories}{left exact} by \Cref{lemma:invariant_coinvariant_functors_for_group_rings_are_adjoint_to_trivial_module_functor}.

    \item The \hldef{homology group of $G$ with coefficients in $A$} or the \hldef{group homology of $A$ as an $R[G]$-module}, denoted \hl{$H_n(G, A)$} for $n \geq 0$, is the $n$-th \CrefAndHyperrefIfExist{definition:left_right_derived_functors_of_a_right_left_exact_functor_between_abelian_categories_where_source_has_enough_projectives_injectives}{left derived functor} object $(L_n(-_G))(A)$ of the \CrefAndHyperrefIfExist{definition:invariant_subgroup_of_a_group_acting_on_an_abelian_group}{coinvariants} functor $A \mapsto A_G$, which is \CrefAndHyperrefIfExist{definition:exact_functor_between_abelian_categories}{right exact} by \Cref{lemma:invariant_coinvariant_functors_for_group_rings_are_adjoint_to_trivial_module_functor}.
\end{enumerate}
 
\end{definition}
\begin{corollary} \label{corollary:degree_zero_group_homology_cohomology_of_a_module_over_a_group_ring_over_a_ring_are_invariants_coinvariants}
    Let $G$ be a \CrefAndHyperrefIfExist{definition:group}{group} and $R$ a \CrefAndHyperrefIfExist{definition:ring}{not necessarily commutative ring}. Write $R[G]$ for the \CrefAndHyperrefIfExist{definition:group_ring_of_a_group_over_a_ring}{group ring} and let $A$ be a \CrefAndHyperrefIfExist{definition:left_right_module_over_a_group_ring_of_a_group_over_a_ring}{(left or right) $R[G]$-module}. We have

    \begin{enumerate}
        \item $H^0(G,A) \cong A^G$ (\Cref{definition:group_cohomology_homology_of_a_module_over_a_group_ring_over_a_ring}) (\Cref{definition:invariant_subgroup_of_a_group_acting_on_an_abelian_group}),
        \item $H_0(G,A) \cong A_G$ (\Cref{definition:group_cohomology_homology_of_a_module_over_a_group_ring_over_a_ring}) (\Cref{definition:coinvariant_quotient_group_of_a_group_acting_on_an_abelian_group}). 
    \end{enumerate}
\end{corollary}
\begin{proof}
    This follows from the definitions of group homology/cohomology and \Cref{lemma:zeroth_left_right_derived_object_of_a_right_left_exact_functor_is_naturally_isomorphic_to_functor_object}.
\end{proof}
\begin{corollary} \label{corollary:group_homology_cohomology_are_just_tor_and_ext_functors_by_trivial_module_over_group_ring}
    Let $G$ be a \CrefAndHyperrefIfExist{definition:group}{group} and $R$ a \CrefAndHyperrefIfExist{definition:ring}{not necessarily commutative ring}. Let $A$ be a \CrefAndHyperrefIfExist{definition:left_right_module_over_a_group_ring_of_a_group_over_a_ring}{left/right $R[G]$-module}.
    \begin{enumerate}
        \item If $A$ is a left $R[G]$-module, then
        $$H_*(G;A) \cong \Tor_*^{R[G]}(R, A)$$
        (\Cref{definition:group_cohomology_homology_of_a_module_over_a_group_ring_over_a_ring}) (\Cref{definition:tor_functors_of_bimodules_of_rings}).
        \item If $A$ is a right $R[G]$-module, then
        $$H_*(G;A) \cong \Tor_*^{R[G]}(A, R)$$
        (\Cref{definition:group_cohomology_homology_of_a_module_over_a_group_ring_over_a_ring}) (\Cref{definition:tor_functors_of_bimodules_of_rings}).

        \item If $A$ is a left $R[G]$-module, then
        $$H^*(G;A) \cong \Ext_{R[G]}^{*}(R, A)$$
        (\Cref{definition:group_cohomology_homology_of_a_module_over_a_group_ring_over_a_ring}) (\Cref{definition:ext_functors_of_bimodules_of_rings}); note here that $\Ext_{R[G]}$ is the \CrefAndHyperrefIfExist{definition:left_right_derived_functors_of_a_right_left_exact_functor_between_abelian_categories_where_source_has_enough_projectives_injectives}{right derived functor} of $\Hom_{R[G]-\bbZ}$ (\Cref{definition:hom_of_left_right_bi_modules_of_rings})
        \item If $A$ is a right $R[G]$-module, then
        $$H^*(G;A) \cong \Ext_{R[G]}^{*}(R, A)$$
        (\Cref{definition:group_cohomology_homology_of_a_module_over_a_group_ring_over_a_ring}) (\Cref{definition:ext_functors_of_bimodules_of_rings}); note here that $\Ext_{R[G]}$ is the \CrefAndHyperrefIfExist{definition:left_right_derived_functors_of_a_right_left_exact_functor_between_abelian_categories_where_source_has_enough_projectives_injectives}{right derived functor} of $\Hom_{\bbZ-R[G]}$ (\Cref{definition:hom_of_left_right_bi_modules_of_rings})

    \end{enumerate}
    In all of the above statements, $R$ is regarded as a \CrefAndHyperrefIfExist{definition:trivial_G_module_of_a_group_over_a_ring}{trivial $R[G]$-module}.
\end{corollary}
\begin{proof}
    By definition, $H_n(G,A)$ is the $n$th left derived functor of the coinvariant functor $-_G$, which is naturally isomorphic to $R \otimes_{R[G]} -$ (for left $R[G]$-modules $A$) by \Cref{lemma:invariant_coinvariant_functors_for_group_rings_are_naturally_isomorphic_to_tensor_hom_by_trivial_module}. The $n$th left derived functor of $R \otimes_{R[G]} -$ is by definition $\Tor_n^{R[G]}(R,-)$. 

    Similarly, argue that $H^*(G;A)$ is the $n$th right derived functor of the invariant functor $-^G$, which is naturally isomorphic to $\Hom_{R[G]}(R, A)$ (for left $R[G]$-modules $A$) by \Cref{lemma:invariant_coinvariant_functors_for_group_rings_are_naturally_isomorphic_to_tensor_hom_by_trivial_module}. The $n$th left derived functor of $\Hom_{R[G]}(R, -)$ is by definition $\Ext^n_{R[G]}(R,-)$. 
\end{proof}

\subsection{Computing group homology/cohomology from Bar resolutions}
\begin{notation} \label{notation:basis_elements_for_n_fold_cartesian_product_of_a_group}
Let $R$ be a \CrefAndHyperrefIfExist{definition:ring}{not necessarily commutative ring} and $G$ a \CrefAndHyperrefIfExist{definition:group}{group}. 
For any $n \geq 0$, let \hl{$G^{n}$} denote the $n$-fold Cartesian product of $G$. An element $(g_1, \dots, g_n) \in G^n$ is denoted by \hl{$[g_1 | \dots | g_n]$}. For $n=0$, we denote the basis element of $G^0$ by \hl{$[]$}.
\end{notation}
\begin{definition} \label{definition:unnormalized_bar_resolution_of_a_ring_over_a_group_algebra}
Let $R$ be a \CrefAndHyperrefIfExist{definition:ring}{not necessarily commutative ring} and $G$ a \CrefAndHyperrefIfExist{definition:group}{group}. Apply \Cref{notation:basis_elements_for_n_fold_cartesian_product_of_a_group}.

The \hldef{unnormalized bar resolution of $R$ over $R[G]$} is the complex \hl{$B_\bullet = B_{\bullet}(R[G])$} of (left) $R[G]$-modules where each $B_n$ is the \CrefAndHyperrefIfExist{definition:free_bimodule_over_rings}{free $R[G]$-module} with \CrefAndHyperrefIfExist{definition:free_bimodule_over_rings}{basis} $G^n$. Noting that the free $R$-module with basis $G^n$ can be described by 
$$R[G]^{\otimes_R n} = \underbrace{R[G] \otimes_R R[G] \otimes_R \cdots \otimes_R R[G]}_{n \text{ times}}$$
where all of the copies of $R[G]$ are considered with the two-sided $R$-module structure, further note that $B_n$, the free $R[G]$-module with basis $G^n$, is thus identifiable as 
$$B_n = R[G] \otimes_R R[G]^{\otimes_R n} \cong \underbrace{R[G] \otimes_R R[G] \dots \otimes_R R[G]}_{n+1 \text{ times}}$$
where the first copy of $R[G]$ is considered with its $R[G]-R$-bimodule structure.

The \hldef{differential} \hl{$d_n: B_n \to B_{n-1}$} is defined on the $R[G]$-basis by:
$$d_n(g_0[g_1 | \dots | g_n]) = g_0 g_1[g_2 | \dots | g_n] + \left( \sum_{i=1}^{n-1} (-1)^i g_0[g_1 | \dots | g_i g_{i+1} | \dots | g_n] \right) + (-1)^n g_0[g_1 | \dots | g_{n-1}]$$
and extended $R[G]$-linearly. 

This resolution is indeed (\Cref{theorem:unnormalized_and_normalized_bar_resolutions_of_group_rings_are_resolutions}) a resolution
$\dots \to B_2 \to B_1 \to B_0 \to R \to 0,$
more precisely a \CrefAndHyperrefIfExist{definition:left_right_resolution_of_a_class_of_objects_in_an_abelian_category}{left resolution} by free modules, of the \CrefAndHyperrefIfExist{definition:trivial_G_module_of_a_group_over_a_ring}{trivial $R[G]$-module} $R$, where the map $\epsilon: B_0 \to R$ is given by $\epsilon(g_0[]) = 1$; note that this is identifiable with the \CrefAndHyperref{definition:augmentation_map_of_a_group_ring_to_its_base_ring}{augmentation map} $\epsilon: R[G] \to R$ given by $\epsilon(g) = 1$ for all $g \in G$.

Symmetrically, one can alternatively construct $B_\bullet$ as a (left free) resolution of $R$ of right $R[G]$-modules.
\end{definition}
\begin{definition} \label{definition:normalized_bar_resolution_of_a_ring_over_a_group_algebra}
Let $R$ be a \CrefAndHyperrefIfExist{definition:ring}{not necessarily commutative ring} and $G$ a \CrefAndHyperrefIfExist{definition:group}{group}. Apply \Cref{notation:basis_elements_for_n_fold_cartesian_product_of_a_group}.

\TODO{subcomplex of a complex, quotient complex of a complex}
The \hldef{normalized bar resolution} \hl{$\overline{B}_\bullet = \overline{B}_{\bullet}(R[G])$} is the \CrefAndHyperrefIfExist{definition:quotient_object_of_an_object_of_an_abelian_category_by_a_subobject}{quotient} of the \CrefAndHyperrefIfExist{definition:unnormalized_bar_resolution_of_a_ring_over_a_group_algebra}{unnormalized bar resolution} \hl{$B_{\bullet}(R[G])$} by the subcomplex \hl{$D_{\bullet}$}, where \hl{$D_n$} is the $R[G]$-\CrefAndHyperrefIfExist{definition:submodule_of_a_module_over_a_ring}{submodule} of \CrefAndHyperrefIfExist{definition:unnormalized_bar_resolution_of_a_ring_over_a_group_algebra}{$B_n$} spanned by elements $g_0[g_1 | \dots | g_n]$ such that at least one $g_i = 1$ for $i \geq 1$. The \CrefAndHyperrefIfExist{definition:module_over_a_sheaf_of_rings_on_a_site}{modules} in the normalized resolution are given by $\overline{B}_n = R[G] \otimes_R (R[G] / R \cdot [1])^{\otimes n}$. Equivalently, $\overline{B}_n$ can be identified as the free $R[G]$-module generated by $(G \setminus \{1\})^n$, i.e. the set of symbols $[g_1 | \dots | g_n]$ where $g_i \in G \setminus \{1\}$.

The \CrefAndHyperrefIfExist{definition:chain_complex_of_objects_in_a_preadditive_category}{differential} is induced by the differential $d_n$ of the unnormalized bar resolution.

This resolution is indeed (\Cref{theorem:unnormalized_and_normalized_bar_resolutions_of_group_rings_are_resolutions}) a resolution
$\dots \to \overline{B}_2 \to \overline{B}_1 \to \overline{B}_0 \to R \to 0,$
more precisely a \CrefAndHyperrefIfExist{definition:left_right_resolution_of_an_object_by_a_class_of_objects_in_an_abelian_category}{left resolution} by free modules, of the \CrefAndHyperrefIfExist{definition:trivial_G_module_of_a_group_over_a_ring}{trivial $R[G]$-module} $R$, where the map $\epsilon: \overline{B}_0 \to R$ is given by $\epsilon(g_0[]) = 1$; note that this is identifiable with the \CrefAndHyperref{definition:augmentation_map_of_a_group_ring_to_its_base_ring}{augmentation map} $\epsilon: R[G] \to R$ given by $\epsilon(g) = 1$ for all $g \in G$.

As with the unnormalized bar resolution, one can construct $\overline{B}_\bullet$ as either a left free resolution of $R$ by left $R[G]$-modules or by right $R[G]$-modules.
\end{definition}
\begin{theorem}[cf {\cite[Theorem 6.5.3]{weibel}}] \label{theorem:unnormalized_and_normalized_bar_resolutions_of_group_rings_are_resolutions}
    Let $R$ be a \CrefAndHyperrefIfExist{definition:ring}{ring} and let $G$ be a \CrefAndHyperrefIfExist{definition:group}{group}.
    The \CrefAndHyperrefIfExist{definition:unnormalized_bar_resolution_of_a_ring_over_a_group_algebra}{unnormalized} and \CrefAndHyperrefIfExist{definition:normalized_bar_resolution_of_a_ring_over_a_group_algebra}{normalized bar resolutions} of $R$ over $R[G]$ are \CrefAndHyperrefIfExist{definition:left_right_resolution_of_an_object_by_a_class_of_objects_in_an_abelian_category}{(left) resolutions} of $R$ by \CrefAndHyperrefIfExist{definition:free_bimodule_over_rings}{free modules} over $R[G]$.  
\end{theorem}



\begin{corollary} \label{corollary:group_homology_cohomology_of_a_module_over_a_group_ring_over_a_ring_can_be_computed_by_bar_resolutions}
Let $G$ be a \CrefAndHyperrefIfExist{definition:group}{group} and $R$ a \CrefAndHyperrefIfExist{definition:ring}{not necessarily commutative ring}. Let $A$ be a \CrefAndHyperrefIfExist{definition:left_right_module_over_a_group_ring_of_a_group_over_a_ring}{left/right $R[G]$-module}.

The \CrefAndHyperrefIfExist{definition:group_cohomology_homology_of_a_module_over_a_group_ring_over_a_ring}{group homology and cohomology} of $G$ with coefficients in $A$ can be computed via the following identifications:

\begin{enumerate}
    \item \textbf{Homology:} 
    \begin{itemize}

        \item If $A$ is a right $R[G]$-module, then $H_n(G, A)$ is the \CrefAndHyperrefIfExist{definition:homology_and_cohomology_objects_for_a_chain_complex_in_an_additive_category}{homologies} of the complexes $A \otimes_{R[G]} B_\bullet$ and $A \otimes_{R[G]} \overline{B}_\bullet$, where $B_\bullet$ and $\overline{B}_\bullet$ are the \CrefAndHyperrefIfExist{definition:unnormalized_bar_resolution_of_a_ring_over_a_group_algebra}{unnormalized} and \CrefAndHyperrefIfExist{definition_normalized_bar_resolution_of_a_ring_over_a_group_algebra}{normalized} bar resolutions constructed as resolutions of left $R[G]$-modules.
        
        \item If $A$ is a left $R[G]$-module, then $H_n(G, A)$ is the homologies of the complexes $B'_\bullet \otimes_{R[G]} A$ and $\overline{B}'_\bullet \otimes_{R[G]} A$, where $B'_\bullet$ and $\overline{B}'_\bullet$ are constructed as resolutions of right $R[G]$-modules.
    
    \end{itemize}

    \item \textbf{Cohomology:}
    
    \begin{itemize}
        \item If $A$ is a left $R[G]$-module, then $H^n(G, A)$ is the \CrefAndHyperrefIfExist{definition:homology_and_cohomology_objects_for_a_chain_complex_in_an_additive_category}{cohomologies} of the complexes $\Hom_{R[G]-\bbZ}(B_\bullet, A)$ and $\Hom_{R[G]-\bbZ}(\overline{B}_\bullet, A)$, where $B_\bullet$ and $\overline{B}_\bullet$ are constructed as resolutions of left $R[G]$-modules.
        
        \item If $A$ is a right $R[G]$-module, then $H^n(G, A)$ is the cohomologies of the complexes $\Hom_{\bbZ-R[G]}(B'_\bullet, A)$ and $\Hom_{\bbZ-R[G]}(\overline{B}'_\bullet, A)$, where $B'_\bullet$ and $\overline{B}'_\bullet$ are constructed as resolutions of right $R[G]$-modules.
    \end{itemize}
\end{enumerate}

% Consequently, these yield the natural isomorphisms to the \CrefAndHyperrefIfExist{definition:tor_functors_of_bimodules_of_rings}{Tor} and \CrefAndHyperrefIfExist{definition:ext_functors_of_bimodules_of_rings}{Ext} functors established in \Cref{corollary:group_homology_cohomology_are_just_tor_and_ext_functors_by_trivial_module_over_group_ring}.
\end{corollary}
\begin{proof}
    This follows from the following:
    \begin{itemize}
        \item \Cref{corollary:group_homology_cohomology_are_just_tor_and_ext_functors_by_trivial_module_over_group_ring}, 
        \item that \CrefAndHyperrefIfExist{definition:tor_functors_of_bimodules_of_rings}{$\Tor$ functors} of modules can be computed with \CrefAndHyperrefIfExist{definition:left_right_resolution_of_a_class_of_objects_in_an_abelian_category}{projective resolutions} in either input, 
        \item that \CrefAndHyperrefIfExist{definition:ext_functors_of_bimodules_of_rings}{$\Ext$ functors} of modules can be computed with projective resolutions in the first input, 
        \item that \CrefAndHyperrefIfExist{definition:free_bimodule_over_rings}{free modules} are \CrefAndHyperrefIfExist{definition:projective_bimodule_over_rings}{projective} (\Cref{proposition:free_bimodule_over_rings_are_projective}).
        \item that the normalized and unnoramlized bar resolutions are (left) resolutions by free $R[G]$-modules (\Cref{theorem:unnormalized_and_normalized_bar_resolutions_of_group_rings_are_resolutions}).
    \end{itemize}
\end{proof}


\subsection{Computatations of group homology/cohomology in simple cases}

\begin{lemma}[cf. {\cite[Lemma 6.1.9]{weibel}}] \label{lemma:invariants_of_group_ring_is_generated_by_norm_element}
    Let $G$ be a \CrefAndHyperrefIfExist{definition:countable_finite_uncountable_sets}{finite} \CrefAndHyperrefIfExist{definition:group}{group}. 
    The \CrefAndHyperrefIfExist{definition:invariant_subgroup_of_a_group_acting_on_an_abelian_group}{invariant subgroup} 
    $$(R[G])^G \cong H^0(G, R[G])$$
    (\Cref{definition:group_cohomology_homology_of_a_module_over_a_group_ring_over_a_ring}) (\Cref{corollary:degree_zero_group_homology_cohomology_of_a_module_over_a_group_ring_over_a_ring_are_invariants_coinvariants})
    of $R[G]$ is the \CrefAndHyperrefIfExist{definition:sheaf_of_ideals_for_a_sheaf_of_rings_on_a_site}{two-sided ideal} of $R[G]$ \CrefAndHyperrefIfExist{definition:ideal_generated_by_a_subset_of_a_ring}{generated by} the \CrefAndHyperrefIfExist{definition:norm_element_of_a_group_algebra_of_a_finite_group_over_a_ring}{norm element} $N$.
\end{lemma}
\begin{proof}
    Clearly, $N \in R[G]^G$. Conversely, if $a = \sum_{g \in G} n_g g$ is in $R[G]^G$, then all of the coefficients $n_g$ are the same. Thus, $a = nN$ for some element $r \in R$. 
\end{proof}

\begin{proposition}[cf. {\cite[Exercise 6.1.3, 1]{weibel}}] \label{proposition:invariants_of_group_ring_for_infinite_group_is_zero}
    Let $G$ be an \CrefAndHyperrefIfExist{definition:countable_finite_uncountable_sets}{infinite} \CrefAndHyperrefIfExist{definition:group}{group}. We have 
    $$H^0(G; R[G]) \cong (R[G])^G = 0.$$
    (\Cref{definition:group_cohomology_homology_of_a_module_over_a_group_ring_over_a_ring}) (\Cref{definition:group_ring_of_a_group_over_a_ring}) (\Cref{definition:invariant_subgroup_of_a_group_acting_on_an_abelian_group}) (\Cref{corollary:degree_zero_group_homology_cohomology_of_a_module_over_a_group_ring_over_a_ring_are_invariants_coinvariants})
\end{proposition}
\begin{proof}
    If $a = \sum_{g \in G} n_g g$ is in $R[G]^G$, then all of the coefficients $n_g$ are the same. However, by definition, elements of $R[G]$ are finite sums, so the $n_g$ must be $0$.
\end{proof}

\subsubsection{Group homology/cohomology of cyclic groups}

\begin{lemma} \label{lemma:group_ring_of_a_cyclic_group_has_periodic_free_resolution}
    Let $m \geq 1$ be an integer. Let $C_m$ be the cyclic \CrefAndHyperrefIfExist{definition:group}{group} of $m$ elements \CrefAndHyperrefIfExist{definition:subgroup_of_a_group_generated_by_elements}{generated by} $\sigma$. Let $R$ be a \CrefAndHyperrefIfExist{definition:ring}{(not necessarily commutative) ring}. Write $N = \sum_{i=0}^{m-1} \sigma^i$ for the \CrefAndHyperrefIfExist{definition:norm_element_of_a_group_algebra_of_a_finite_group_over_a_ring}{norm element}.

    The following is a \CrefAndHyperrefIfExist{definition:left_right_resolution_of_an_object_by_a_class_of_objects_in_an_abelian_category}{left resolution} of the \CrefAndHyperrefIfExist{definition:trivial_G_module_of_a_group_over_a_ring}{trivial $R[C_m]$-module} $R$ by \CrefAndHyperrefIfExist{definition:free_bimodule_over_rings}{free} $R[C_m]$-\CrefAndHyperrefIfExist{definition:module_of_a_cosmos_between_categories_enriched_in_the_cosmos}{modules}:
    $$\cdots \xrightarrow{N} R[C_m] \xrightarrow{\sigma-1} R[C_m] \xrightarrow{N} R[C_m] \xrightarrow{\sigma-1} R[C_m] \xrightarrow{\varepsilon} R \to 0$$
    (\Cref{definition:augmentation_map_of_a_group_ring_to_its_base_ring})
\end{lemma}
\input{../_concepts/proposition_group_homology_cohomology_for_cyclic_group.tex}

\subsubsection{First group homology of trivial action on base ring}

\begin{lemma}[{\cite[Exercise 6.1.4]{weibel}}] \label{lemma:quotient_of_augmentation_ideal_of_group_by_square_is_isomorphic_to_abelianization_of_group}
    Let $G$ be a \CrefAndHyperrefIfExist{definition:group}{group}. There is an isomorphism
    $$I_G / I_G^2 \cong G/[G,G]$$
    \TODO{commutator}
    of abelian groups where $I_G$ is the \CrefAndHyperrefIfExist{definition:augmentation_map_of_a_group_ring_to_its_base_ring}{augmentation ideal} of $\bbZ[G]$ and $[G,G]$ is the commutator subgroup of $G$. 
\end{lemma}
\begin{proposition}[{\cite[Theorem 6.1.11]{weibel}}]  \label{proposition:first_group_homology_of_a_trivial_G_module_is_abelianization_of_group_tensored_with_module}
    Let $G$ be a \CrefAndHyperrefIfExist{definition:group}{group}. Let $M$ be an \CrefAndHyperrefIfExist{definition:group}{abelian group} \CrefAndHyperrefIfExist{lemma:posets_correspond_to_small_filtered_thin_categories}{regarded} as a \CrefAndHyperrefIfExist{definition:trivial_G_module_of_a_group_over_a_ring}{trivial $G$-module}.
    We have 
    $$H_1(G,M) \cong (G/[G,G]) \otimes_{\bbZ} M$$
    (\Cref{definition:group_cohomology_homology_of_a_module_over_a_group_ring_over_a_ring}) (\Cref{definition:tensor_product_of_bimodules_of_rings})
    \TODO{commutator}
\end{proposition}
\begin{proof}
    \TODO{}
\end{proof}
\begin{proposition}[cf. {\cite[Exercise 6.1.5]{weibel}}] \label{proposition:first_group_cohomology_of_a_module_of_a_group_ring_over_a_commutative_ring_is_isomorphic_to_homs}
    Let $G$ be a \CrefAndHyperrefIfExist{definition:group}{group}. Let $R$ be a \CrefAndHyperrefIfExist{definition:ring}{commutative ring}. Let $A$ be a \CrefAndHyperrefIfExist{definition:trivial_G_module_of_a_group_over_a_ring}{trivial $R[G]$-module}. Let $I$ be the \CrefAndHyperrefIfExist{definition:augmentation_map_of_a_group_ring_to_its_base_ring}{augmentation ideal} of $R[G]$.

    There is a natural isomorphism
    $$H^1(G,A) \cong \Hom_{R-\Mod}(I/I^2, A) \cong \Hom_{R-\Mod}(R \otimes_{\bbZ} G/[G,G], A) \cong \Hom_{\mathbf{Grp}}(G,A) \cong \Hom_{\mathbf{Ab}}(G/[G,G], A).$$
    (\Cref{definition:group_cohomology_homology_of_a_module_over_a_group_ring_over_a_ring})\CrefIfExists{definition:category_of_modules_and_bimodules_over_rings}\CrefIfExists{definition:category_of_groups_of_abelian_groups}
    \TODO{commutator}
\end{proposition}




\TODO{e.g. weibel theorem 6.2.2, but do it over general rings}
\TODO{define Tate cohomology group weibel definition 6.2.4}

\subsection{Induction/Coinduction/Restriction of $G$-modules}


\input{../_definitions/definition_restriction_of_scalars_of_a_left_or_right_module_along_a_ring_homomorphism_from_the_base_ring.tex}
\input{../_definitions/definition_extension_of_scalars_base_change_induction_for_modules_along_ring_homomorphisms.tex}
\input{../_definitions/definition_coextension_coinduction_of_scalars_of_modules_along_a_ring_homomorphism.tex}
% \begin{theorem}[Extension-Restriction and Restriction-Coextension adjunction for modules] \label{theorem:extension_restriction_and_restriction_coextension_adjunction_for_modules_over_rings}
% Let $R$ and $S$ be \CrefAndHyperrefIfExist{definition:ring}{rings} (not necessarily commutative), and let $f: R \to S$ be a \CrefAndHyperrefIfExist{definition:ring_homomorphism}{ring homomorphism}.
%
% \begin{enumerate}
%     \item \textbf{Extension-Restriction adjunction:} The \CrefAndHyperrefIfExist{definition:extension_of_scalars_base_change_induction_for_modules_along_ring_homomorphisms}{extension of scalars functor} $f^* \coloneqq S \otimes_R -$ is \CrefAndHyperrefIfExist{definition:adjoint_functors_between_categories_unit_counit_of_adjoint_functors}{left adjoint} to the \CrefAndHyperrefIfExist{definition:restriction_of_scalars_of_a_left_or_right_module_along_a_ring_homomorphism_from_the_base_ring}{restriction of scalars functor} $f_*$.
%
%     Explicitly, for any left $R$-module $M$ and any left $S$-module $N$, there is a \CrefAndHyperrefIfExist{definition:natural_transformation_between_functors_between_categories}{natural isomorphism} of abelian groups:
%     \[
%         \Hom_S(S \otimes_R M, N) \cong \Hom_R(M, f_*N)
%     \]
%     given by the mapping
%     \[
%         \Psi \mapsto (m \mapsto \Psi(1 \otimes m)).
%     \]
%
%     \item \textbf{Restriction-Coextension adjunction:} The \CrefAndHyperrefIfExist{definition:coextension_coinduction_of_scalars_of_modules_along_a_ring_homomorphism}{co-extension of scalars functor} $f^! \coloneqq \Hom_R(S, -)$ is \CrefAndHyperrefIfExist{definition:adjoint_functors_between_categories_unit_counit_of_adjoint_functors}{right adjoint} to the restriction of scalars functor $f_*$.
%
%     Explicitly, for any left $S$-module $N$ and any left $R$-module $M$, there is a \CrefAndHyperrefIfExist{definition:natural_transformation_between_functors_between_categories}{natural isomorphism} of abelian groups:
%     \[ \Hom_R(f_*N, M) \cong \Hom_S(N, \Hom_R(S, M)) \]
%     where $S$ is viewed as an $(R, S)$-bimodule when constructing the the \CrefAndHyperrefIfExist{definition:hom_of_left_right_bi_modules_of_rings}{Hom-module $\Hom_R(S, M)$}.
% \end{enumerate}
% \end{theorem}
%
\begin{theorem}[Extension-Restriction and Restriction-Coextension adjunction for modules] \label{theorem:extension_restriction_and_restriction_coextension_adjunction_for_modules_over_rings}
Let $R$ and $S$ be \CrefAndHyperrefIfExist{definition:ring}{rings} (not necessarily commutative), and let $\varphi : R \to S$ be a \CrefAndHyperrefIfExist{definition:ring_homomorphism}{ring homomorphism}. Let $T$ be a ring.

\begin{enumerate}
    \item \textbf{Extension-Restriction adjunction:}
    \begin{enumerate}
        \item \textbf{Restriction on the Left:} The \CrefAndHyperrefIfExist{definition:extension_of_scalars_base_change_induction_for_modules_along_ring_homomorphisms}{extension of scalars functor}
        \[ f^* = S \otimes_R - : {}_R\mathsf{Mod}_T \to {}_S\mathsf{Mod}_T \]
        \CrefIfExists{definition:category_of_modules_and_bimodules_over_rings}
        is \CrefAndHyperrefIfExist{definition:adjoint_functors_between_categories_unit_counit_of_adjoint_functors}{left adjoint} to the \CrefAndHyperrefIfExist{definition:restriction_of_scalars_of_a_left_or_right_module_along_a_ring_homomorphism_from_the_base_ring}{restriction of scalars functor}
        \[ f_* : {}_S\mathsf{Mod}_T \to {}_R\mathsf{Mod}_T. \]
        Explicitly, for any $R$-$T$-bimodule $M$ and any $S$-$T$-bimodule $N$, there is a natural isomorphism:
        \[
            \Hom_{S-T}(S \otimes_R M, N) \cong \Hom_{R-T}(M, f_*N).
        \]

        \item \textbf{Restriction on the Right:} The extension of scalars functor
        \[ f^* = - \otimes_R S : {}_T\mathsf{Mod}_R \to {}_T\mathsf{Mod}_S \]
        is left adjoint to the restriction of scalars functor
        \[ f_* : {}_T\mathsf{Mod}_S \to {}_T\mathsf{Mod}_R. \]
        Explicitly, for any $T$-$R$-bimodule $M$ and any $T$-$S$-bimodule $N$, there is a natural isomorphism:
        \[
            \Hom_{T-S}(M \otimes_R S, N) \cong \Hom_{T-R}(M, f_*N).
        \]
    \end{enumerate}

    \item \textbf{Restriction-Coextension adjunction:}
    \begin{enumerate}
        \item \textbf{Restriction on the Left:} The \CrefAndHyperrefIfExist{definition:coextension_coinduction_of_scalars_of_modules_along_a_ring_homomorphism}{co-extension of scalars functor}
        \[ f^! = \Hom_R(S, -) : {}_R\mathsf{Mod}_T \to {}_S\mathsf{Mod}_T \]
        is \CrefAndHyperrefIfExist{definition:adjoint_functors_between_categories_unit_counit_of_adjoint_functors}{right adjoint} to the restriction of scalars functor
        \[ f_* : {}_S\mathsf{Mod}_T \to {}_R\mathsf{Mod}_T. \]
        Explicitly, for any $S$-$T$-bimodule $N$ and any $R$-$T$-bimodule $M$, there is a natural isomorphism:
        \[ \Hom_{R-T}(f_*N, M) \cong \Hom_{S-T}(N, \Hom_R(S, M)). \]

        \item \textbf{Restriction on the Right:} The co-extension of scalars functor
        \[ f^! = \Hom_R(S, -) : {}_T\mathsf{Mod}_R \to {}_T\mathsf{Mod}_S \]
        is right adjoint to the restriction of scalars functor
        \[ f_* : {}_T\mathsf{Mod}_S \to {}_T\mathsf{Mod}_R. \]
        Explicitly, for any $T$-$S$-bimodule $N$ and any $T$-$R$-bimodule $M$, there is a natural isomorphism:
        \[ \Hom_{T-R}(f_*N, M) \cong \Hom_{T-S}(N, \Hom_R(S, M)) \]
        where here $\Hom_R(S, M)$ uses the left $R$-module structure on $S$ and the right $R$-module structure on $M$.
    \end{enumerate}
\end{enumerate}
\end{theorem}

\begin{proof}
    \TODO{AI generated, verify}
    We prove part (1)(a), the left-restriction extension-Restriction adjunction. The other three cases follow by identical arguments with left/right roles swapped.

    Let $M$ be an $R$-$T$ bimodule and $N$ be an $S$-$T$ bimodule. For any $S$-$T$ bimodule homomorphism $\Psi : S \otimes_R M \to N$, define
    $$\Phi(\Psi) : M \to f_*N, \quad m \mapsto \Psi(1_S \otimes m).$$
    We verify that $\Phi(\Psi)$ is an $R$-$T$ bimodule homomorphism. For any $r \in R$, $t \in T$, and $m \in M$:
    $$\Phi(\Psi)(r \cdot m) = \Psi(1_S \otimes r \cdot m).$$
    In the tensor product $S \otimes_R M$, we have $1_S \otimes r \cdot m = 1_S \otimes \varphi(r)m = \varphi(r) \cdot 1_S \otimes m = (\varphi(r) \cdot 1_S) \otimes m$. Since $\Psi$ is an $S$-$T$ bimodule homomorphism and hence in particular $R$-linear (via restriction), we obtain
    $$\Phi(\Psi)(r \cdot m) = \Psi(\varphi(r) \cdot (1_S \otimes m)) = \varphi(r) \cdot \Psi(1_S \otimes m).$$
    By the definition of restriction of scalars\CrefIfExists{definition:restriction_of_scalars_of_a_left_or_right_module_along_a_ring_homomorphism_from_the_base_ring}, we have $r \cdot_R x = \varphi(r) \cdot_S x$ for any element $x$ of an $S$-module. Thus, $\Phi(\Psi)(r \cdot m) = r \cdot_R \Phi(\Psi)(m)$, and similarly $\Phi(\Psi)(m \cdot t) = \Phi(\Psi)(m) \cdot t$. Therefore, $\Phi(\Psi) \in \operatorname{Hom}_{R-T}(M, f_*N)$.

    Next we construct the inverse map. Given an $R$-$T$ bimodule homomorphism $g : M \to f_*N$, consider the map
    $$\gamma_g : S \times M \to N, \quad (s, m) \mapsto s \cdot g(m).$$
    We verify that $\gamma_g$ satisfies the hypotheses of \Cref{proposition:universal_property_of_the_tensor_product_of_bimodules}. For all $r \in R$, $t \in T$, $m,m' \in M$, and $s \in S$:
    $$\gamma_g(s, m+m') = s \cdot g(m+m') = s \cdot (g(m) + g(m')) = \gamma_g(s,m) + \gamma_g(s,m'),$$
    where the second equality uses that $g$ is additive and the third uses that scalar multiplication in $N$ distributes over addition. Similarly, $\gamma_g(s+s', m) = (s+s') \cdot g(m) = s \cdot g(m) + s' \cdot g(m) = \gamma_g(s,m) + \gamma_g(s',m)$, and
    $$\gamma_g(s, m \cdot t) = s \cdot g(m \cdot t) = s \cdot (g(m) \cdot t) = (s \cdot g(m)) \cdot t = \gamma_g(s,m) \cdot t,$$
    where the second equality uses that $g$ is a $T$-bimodule homomorphism and the third uses the bimodule axioms on $N$. Finally, for the balancing condition over $R$:
    $$\gamma_g(s \cdot \varphi(r), m) = (s \cdot \varphi(r)) \cdot g(m) = s \cdot (\varphi(r) \cdot g(m)) = s \cdot (r \cdot_R g(m)),$$
    while also
    $$\gamma_g(s, r \cdot m) = s \cdot g(r \cdot m).$$
    Since $g$ is an $R$-$T$ bimodule homomorphism from $M$ to $f_*N$, we have $g(r \cdot m) = r \cdot_R g(m) = \varphi(r) \cdot_S g(m)$ by the definition of restriction of scalars. Therefore,
    $$\gamma_g(s, r \cdot m) = s \cdot (\varphi(r) \cdot g(m)) = (s \cdot \varphi(r)) \cdot g(m) = \gamma_g(\varphi(r)s, m).$$
    This is the correct balancing condition for the tensor product $S \otimes_R M$, where $S$ is viewed as an $(S,R)$-bimodule via $\varphi$. By \Cref{proposition:universal_property_of_the_tensor_product_of_bimodules}, there exists a unique group homomorphism
    $$\Theta(g) : S \otimes_R M \to N$$
    such that $\Theta(g)(s \otimes m) = s \cdot g(m)$ for all $s \in S$, $m \in M$.

    We verify that $\Theta(g)$ is an $S$-$T$ bimodule homomorphism. For any $s' \in S$ and pure tensor $s \otimes m$:
    $$\Theta(g)(s' \cdot (s \otimes m)) = \Theta(g)((s's) \otimes m) = (s's) \cdot g(m) = s' \cdot (s \cdot g(m)) = s' \cdot \Theta(g)(s \otimes m).$$
    For the right $T$-action:
    $$\Theta(g)((s \otimes m) \cdot t) = \Theta(g)(s \otimes (m \cdot t)) = s \cdot g(m \cdot t) = s \cdot (g(m) \cdot t) = (s \cdot g(m)) \cdot t = \Theta(g)(s \otimes m) \cdot t.$$
    Since both actions are additive, these equalities extend to all of $S \otimes_R M$. Thus, $\Theta(g) \in \operatorname{Hom}_{S-T}(S \otimes_R M, N)$.

    We verify that $\Phi$ and $\Theta$ are mutual inverses. For any $\Psi : S \otimes_R M \to N$:
    $$\Theta(\Phi(\Psi))(1_S \otimes m) = 1_S \cdot \Phi(\Psi)(m) = 1_S \cdot \Psi(1_S \otimes m) = \Psi(1_S \otimes m).$$
    Since both $\Theta(\Phi(\Psi))$ and $\Psi$ are $S$-$T$ bimodule homomorphisms agreeing on all pure tensors of the form $1_S \otimes m$, and every element of $S \otimes_R M$ is a finite sum of elements $s \otimes m = s(1_S \otimes m)$, we have $\Theta(\Phi(\Psi)) = \Psi$. Conversely, for any $g : M \to f_*N$:
    $$\Phi(\Theta(g))(m) = \Theta(g)(1_S \otimes m) = 1_S \cdot g(m) = g(m).$$
    Therefore, $\Phi(\Theta(g)) = g$, and the correspondence is bijective.

    For the Restriction-Coextension adjunction (part (2)(a)), let $N$ be an $S$-$T$ bimodule and $M$ be an $R$-$T$ bimodule. Given an $R$-$T$ bimodule homomorphism $g : f_*N \to M$, define
    $$\Psi_g : N \to \Hom_R(S, M), \quad n \mapsto (s \mapsto g(s \cdot n)).$$
    We verify that for each $n \in N$, the map $\Psi_g(n) : S \to M$ is a left $R$-module homomorphism. For all $r \in R$ and $s \in S$:
    $$\Psi_g(n)(rs) = g((rs) \cdot n).$$
    The element $rs$ acts on $n$ via the $(S,R)$-bimodule structure: $rs \cdot n = r \cdot_S s \cdot n = \varphi(r)s \cdot n$. Since $g$ is an $R$-$T$ bimodule homomorphism, we have
    $$\Psi_g(n)(rs) = g(\varphi(r)s \cdot n) = r \cdot g(s \cdot n) = r \cdot \Psi_g(n)(s),$$
    so $\Psi_g(n) \in \operatorname{Hom}_R(S, M)$ by the definition of left $R$-linearity.

    We verify that $\Psi_g : N \to \Hom_R(S, M)$ is an $S$-$T$ bimodule homomorphism. For any $s' \in S$, $t \in T$, and $n \in N$:
    $$\bigl(\Psi_g(s' \cdot n)\bigr)(s) = g(s \cdot s' \cdot n) = g((ss') \cdot n).$$
    On the other hand, by the definition of the left $S$-action on $\Hom_R(S,M)$ from \Cref{definition:hom_of_left_right_bi_modules_of_rings}:
    $$(s' \cdot \Psi_g(n))(s) = \Psi_g(n)(s \cdot s') = g((ss') \cdot n).$$
    Thus, $\Psi_g(s' \cdot n) = s' \cdot \Psi_g(n)$, and similarly for the right $T$-action. Therefore, $\Psi_g \in \operatorname{Hom}_{S-T}(N, \Hom_R(S, M))$.

    We construct the inverse. Given $h : N \to \Hom_R(S, M)$ in $\operatorname{Hom}_{S-T}(N, \Hom_R(S, M))$, define
    $$\Gamma(h) : f_*N \to M, \quad n \mapsto h(n)(1_S).$$
    We verify that $\Gamma(h)$ is an $R$-$T$ bimodule homomorphism. For any $r \in R$:
    $$\Gamma(h)(r \cdot_R n) = \Gamma(h)(\varphi(r) \cdot n) = h(\varphi(r) \cdot n)(1_S).$$
    Since $h$ is an $S$-$T$ bimodule homomorphism, we have $h(\varphi(r) \cdot n) = \varphi(r) \cdot h(n)$ in $\Hom_R(S, M)$. By the left $S$-action on $\Hom_R(S, M)$:
    $$\bigl(\varphi(r) \cdot h(n)\bigr)(1_S) = h(n)(1_S \cdot \varphi(r)) = h(n)(\varphi(r)).$$
    On the other hand, since $h(n) : S \to M$ is a left $R$-module homomorphism:
    $$h(n)(\varphi(r)) = h(n)(r \cdot_S 1_S) = r \cdot h(n)(1_S).$$
    Therefore, $\Gamma(h)(r \cdot_R n) = r \cdot h(n)(1_S) = r \cdot \Gamma(h)(n)$, as required. Similarly for the right $T$-action.

    We verify that $\Psi$ and $\Gamma$ are mutual inverses. For any $g : f_*N \to M$:
    $$\bigl(\Psi_{\Gamma(g)}(n)\bigr)(s) = \Gamma(g)(s \cdot n) = g(s \cdot n),$$
    so $\Psi_{\Gamma(g)}(n) = g(s \cdot n)$ for all $s$, which is exactly the definition of $\Psi_g(n)$. Thus, $\Psi_{\Gamma(g)} = g$. Conversely, for any $h : N \to \Hom_R(S, M)$:
    $$\Gamma(\Psi_h)(n) = \Psi_h(n)(1_S) = h(n)(1_S \cdot 1_S) = h(n)(1_S).$$
    Therefore, $\Gamma(\Psi_h) = h$. The correspondence is bijective and natural in both arguments by construction.

    Cases (1)(b), (2)(b) follow by the same constructions with left/right roles interchanged.
\end{proof}



\begin{definition} \label{definition:restriction_of_a_group_module_over_a_ring_to_a_subgroup}
    Let $R$ be a \CrefAndHyperrefIfExist{definition:ring}{not necessarily commutative ring}, let $G$ be a \CrefAndHyperrefIfExist{definition:group}{group}, and let $H \leq G$ be a \CrefAndHyperrefIfExist{definition:subgroup_of_a_group}{subgroup}. Write $i: R[H] \hookrightarrow R[G]$ for the natural inclusion of group algebras.

    Given a \CrefAndHyperrefIfExist{definition:left_right_module_over_a_group_ring_of_a_group_over_a_ring}{$R[G]$-module} $V$, the \hldef{restriction of $V$ to $H$}, denoted \hl{$\text{Res}^G_H(V)$}, is the $H$-module obtained by pulling back the action of $R[G]$ along $i$. As an \CrefAndHyperrefIfExist{definition:bimodule_over_rings_module_over_a_ring}{$R$-module}, $\text{Res}^G_H(V) = V$, where the action is defined by $h \cdot v = i(h)v$ for all $h \in H$ and $v \in V$. 

    In the context of \CrefAndHyperrefIfExist{definition:ring}{ring} theory, this is the \CrefAndHyperrefIfExist{definition:restriction_of_scalars_of_a_left_or_right_module_along_a_ring_homomorphism_from_the_base_ring}{restriction of scalars}.
\end{definition}
\begin{definition} \label{definition:induction_of_a_group_module_over_a_ring_to_a_supergroup}
    Let $R$ be a \CrefAndHyperrefIfExist{definition:ring}{not necessarily commutative ring}, let $G$ be a \CrefAndHyperrefIfExist{definition:group}{group}, and let $H \leq G$ be a \CrefAndHyperrefIfExist{definition:subgroup_of_a_group}{subgroup}. Write $i: R[H] \hookrightarrow R[G]$ for the natural inclusion of group algebras.

    Given an \CrefAndHyperrefIfExist{definition:left_right_module_over_a_group_ring_of_a_group_over_a_ring}{(left) $R[H]$-module} $W$, the \hldef{induction of $W$ to $G$}, denoted \hl{$\text{Ind}^G_H(W)$}, is the $R[G]$-module defined by:
    $$\text{Ind}^G_H(W) = R[G] \otimes_{R[H]} W$$
    where $R[G]$ is viewed as an $(R[G], R[H])$-bimodule.

    In the context of \CrefAndHyperrefIfExist{definition:ring}{ring} theory, this is the \CrefAndHyperrefIfExist{definition:extension_of_scalars_base_change_induction_for_modules_along_ring_homomorphisms}{extension of scalars} of the $R[H]$-module $W$ to $R[G]$.
\end{definition}
\begin{definition} \label{definition:coinduction_of_a_group_module_over_a_ring_to_a_supergroup}
    Let $R$ be a \CrefAndHyperrefIfExist{definition:ring}{not necessarily commutative ring}, let $G$ be a \CrefAndHyperrefIfExist{definition:group}{group}, and let $H \leq G$ be a \CrefAndHyperrefIfExist{definition:subgroup_of_a_group}{subgroup}. Write $i: R[H] \hookrightarrow R[G]$ for the natural inclusion of group algebras.

    Given an (left) $R[H]$-module $W$, the \hldef{coinduction of $W$ to $G$}, denoted \hl{$\text{Coind}^G_H(W)$}, is the $G$-module defined by
    $$\text{Coind}^G_H(W) = \text{Hom}_{R[H]}(R[G], W)$$
    where the $G$-action is defined by $(g \cdot \phi)(x) = \phi(xg)$ for $g, x \in G$ and $\phi \in \text{Hom}_{R[H]}(R[G], W)$. 

    In the context of \CrefAndHyperrefIfExist{definition:ring}{ring} theory, this is the \CrefAndHyperrefIfExist{definition:coextension_coinduction_of_scalars_of_modules_along_a_ring_homomorphism}{coextension of scalars} of the $R[H]$-module $W$ to $R[G]$.
\end{definition}

\begin{proposition} \label{proposition:induction_of_a_group_module_over_a_ring_to_a_supergroup_embeds_into_the_coinduction_and_is_an_isomorphism_if_the_index_is_finite}
    Let $R$ be a \CrefAndHyperrefIfExist{definition:ring}{not necessarily commutative ring}, let $G$ be a \CrefAndHyperrefIfExist{definition:group}{group}, and let $H \leq G$ be a \CrefAndHyperrefIfExist{definition:subgroup_of_a_group}{subgroup}. Write $i: R[H] \hookrightarrow R[G]$ for the natural inclusion of group algebras.

    There is a natural embedding of $G$-modules $\text{Ind}^G_H(W) \to \text{Coind}^G_H(W)$ (\Cref{definition:induction_of_a_group_module_over_a_ring_to_a_supergroup}) (\Cref{definition:coinduction_of_a_group_module_over_a_ring_to_a_supergroup}). If the \CrefAndHyperrefIfExist{definition:index_of_a_subgroup_in_a_group}{index} $[G:H]$ is finite, this map is an isomorphism:
    $$\text{Ind}^G_H(W) \cong \text{Coind}^G_H(W)$$
\end{proposition}
\begin{theorem}[Frobenius reciprocity] \label{theorem:induction_restriction_restriction_coinduction_adjunction_for_group_algebras_over_rings}
    Let $R$ be a \CrefAndHyperrefIfExist{definition:ring}{not necessarily commutative ring}, let $G$ be a \CrefAndHyperrefIfExist{definition:group}{group}, and let $H \leq G$ be a \CrefAndHyperrefIfExist{definition:subgroup_of_a_group}{subgroup}. 

    The \CrefAndHyperrefIfExist{definition:restriction_of_a_group_module_over_a_ring_to_a_subgroup}{restriction functor $\text{Res}^G_H$} admits both a \CrefAndHyperrefIfExist{definition:adjoint_functors_between_categories_unit_counit_of_adjoint_functors}{left adjoint} and a right adjoint:
    \begin{enumerate}
        \item \CrefAndHyperrefIfExist{definition:induction_of_a_group_module_over_a_ring_to_a_supergroup}{Induction} is the left adjoint to restriction:
        $$\text{Hom}_{R[G]}(\text{Ind}^G_H(W), V) \cong \text{Hom}_{R[H]}(W, \text{Res}^G_H(V))$$
        \item Restriction is the left adjoint to \CrefAndHyperrefIfExist{definition:coinduction_of_a_group_module_over_a_ring_to_a_supergroup}{coinduction}:
        $$\text{Hom}_{R[H]}(\text{Res}^G_H(V), W) \cong \text{Hom}_{R[G]}(V, \text{Coind}^G_H(W))$$
    \end{enumerate}
\end{theorem}
\begin{proof}
    This is an immediate reformulation of \Cref{theorem:extension_restriction_and_restriction_coextension_adjunction_for_modules_over_rings} in the case that the ring homomorphism is the inclusion $R[H] \hookrightarrow R[G]$.
\end{proof}

\begin{theorem}[Shapiro's lemma] \label{theorem:shapiros_lemma_for_group_homology_cohomology_of_a_module_over_a_group_ring_over_a_ring_and_of_its_induced_or_coinduced_module}
    Let $R$ be a \CrefAndHyperrefIfExist{definition:ring}{not necessarily commutative ring}, let $G$ be a \CrefAndHyperrefIfExist{definition:group}{group}, and let $H \leq G$ be a \CrefAndHyperrefIfExist{definition:subgroup_of_a_group}{subgroup}. 

    There are \CrefAndHyperrefIfExist{definition:natural_transformation_between_functors_between_categories}{natural isomorphisms} in \CrefAndHyperrefIfExist{definition:group_cohomology_homology_of_a_module_over_a_group_ring_over_a_ring}{group cohomology and homology} natural in \CrefAndHyperrefIfExist{definition:left_right_module_over_a_group_ring_of_a_group_over_a_ring}{$R[H]$-modules} $W$:
    \begin{enumerate}
        \item $H^n(G, \text{Coind}^G_H(W)) \cong H^n(H, W)$ for all $n \geq 0$.
        \item $H_n(G, \text{Ind}^G_H(W)) \cong H_n(H, W)$ for all $n \geq 0$.
    \end{enumerate}
\end{theorem}

\begin{proof}
    Let $P_\bullet \to \mathbb{Z}$ be a projective $\mathbb{Z}G$-resolution. Since $\mathbb{Z}G$ is a free $\mathbb{Z}H$-module, $P_\bullet$ is also a projective $\mathbb{Z}H$-resolution.
    For homology:
    $$ P \otimes_{\mathbb{Z}G} (\mathbb{Z}G \otimes_{\mathbb{Z}H} A) \cong (P \otimes_{\mathbb{Z}G} \mathbb{Z}G) \otimes_{\mathbb{Z}H} A \cong P \otimes_{\mathbb{Z}H} A $$
    The left side computes $H_*(G, \Ind_H^G A)$ and the right side computes $H_*(H, A)$.
    For cohomology, we use the Hom-tensor adjunction:
    $$ \Hom_G(P, \Hom_H(\mathbb{Z}G, A)) \cong \Hom_H(P \otimes_{\mathbb{Z}G} \mathbb{Z}G, A) \cong \Hom_H(P, A) $$
    taking cohomology on both sides yields the result.
\end{proof}

\begin{theorem}[Hilbert's theorem 90, additive version, see e.g. {\cite[6.3.7]{weibel}}] \label{theorem:hliberts_theorem_90_additive_version}
    Let $L/K$ be a \CrefAndHyperrefIfExist{definition:finite_field_extension}{finite} \CrefAndHyperrefIfExist{definition:galois_extension_of_fields}{Galois extension of fields} with \CrefAndHyperrefIfExist{definition:galois_extension_of_fields}{Galois group} $G$. Note that $L$ is a \CrefAndHyperrefIfExist{definition:left_right_module_over_a_group_ring_of_a_group_over_a_ring}{$K[G]$-module} with $L^G, L_G \cong K$\CrefIfExists{lemma:invariant_coinvariant_functors_for_group_rings_are_adjoint_to_trivial_module_functor} as $K$-vector spaces. 
    We have 
    $$H^n(G,L), H_n(G,L) = 0 \text{ for } n \neq 0.$$
    \CrefIfExists{definition:group_cohomology_homology_of_a_module_over_a_group_ring_over_a_ring}
\end{theorem}
\begin{proof}
    By the normal basis theorem \TODO{normal basis theorem}, there is an $x \in L$ such that $\{g(x): g\in G\}$ is a \CrefAndHyperrefIfExist{definition:basis_of_a_vector_space_over_a_field}{basis} of $L$ as a \CrefAndHyperrefIfExist{definition:vector_space_over_a_field}{$K$-vector space}. Thus, $L \cong K[G] \cong \mathbb{Z}[G] \otimes_{\mathbb{Z}} K = \Ind_{\{1\}}^G K$ as $K[G]$-modules. By \CrefAndHyperrefIfExist{theorem:shapiros_lemma_for_group_homology_cohomology_of_a_module_over_a_group_ring_over_a_ring_and_of_its_induced_or_coinduced_module}{Shapiro's lemma}, 
    $$H_n(G, L) \cong H_n(G, \Ind_{\{1\}}^G K) \cong H_n(\{1\}, K) = 0.$$
    Moreover, recall by \Cref{proposition:induction_of_a_group_module_over_a_ring_to_a_supergroup_embeds_into_the_coinduction_and_is_an_isomorphism_if_the_index_is_finite} that $\Ind_{\{1\}}^G(K) \cong \CoInd_{\{1\}}^G(K)$ so we have 
    $$H^n(G, L) \cong H^n(G, \CoInd_{\{1\}}^G K) \cong H^n(\{1\}, K) = 0$$
    by Shapiro's lemma as well.
\end{proof}

\begin{theorem}[Hilbert's theorem 90, multiplicative version, see e.g. {\cite[6.3.10]{weibel}}] \label{theorem:hilberts_theorem_90_multiplicative_version}
    Let $L/K$ be a \CrefAndHyperrefIfExist{definition:degree_of_a_field_extension}{finite} \CrefAndHyperrefIfExist{definition:galois_extension_of_fields}{Galois extension of fields} with \CrefAndHyperrefIfExist{definition:galois_extension_of_fields}{Galois group} $G$. Note that the multiplicative group $L^\times$ is a \CrefAndHyperrefIfExist{definition:left_right_module_over_a_group_ring_of_a_group_over_a_ring}{$G$-module} with $(L^\times)^G = K^\times$. 
    We have 
    $$H^1(G, L^\times) = 0.$$
    \CrefIfExists{definition:group_cohomology_homology_of_a_module_over_a_group_ring_over_a_ring}
\end{theorem}

\begin{proof}
    \TODO{verify}
    Let $\{a_\sigma\}_{\sigma \in G}$ be a \CrefAndHyperrefIfExist{definition:1_cocycle}{1-cocycle} in $Z^1(G, L^\times)$, so $a_{\sigma\tau} = a_\sigma \sigma(a_\tau)$ for all $\sigma, \tau \in G$. By the \CrefAndHyperrefIfExist{theorem:dedekinds_lemma_linear_independence_of_characters}{linear independence of characters}, the map $b: L \to L$ defined by $b(z) = \sum_{\tau \in G} a_\tau \tau(z)$ is not the zero map. Thus, there exists some $z \in L$ such that $\alpha := b(z) \in L^\times$. For any $\sigma \in G$, we calculate:
    $$\sigma(\alpha) = \sum_{\tau \in G} \sigma(a_\tau) \sigma\tau(z)$$
    Using the cocycle identity $\sigma(a_\tau) = a_\sigma^{-1} a_{\sigma\tau}$, we obtain:
    $$\sigma(\alpha) = a_\sigma^{-1} \sum_{\tau \in G} a_{\sigma\tau} \sigma\tau(z) = a_\sigma^{-1} \alpha$$
    Rearranging yields $a_\sigma = \alpha \sigma(\alpha)^{-1}$, which shows that the cocycle is a \CrefAndHyperrefIfExist{definition:1_coboundary}{1-coboundary} in $B^1(G, L^\times)$. Hence, $H^1(G, L^\times) = 0$.
\end{proof}

\subsection{Modules over topological groups and topological rings}

One general class of topological groups is that of profinite groups

\input{../_definitions/definition_profinite_group.tex}
\input{../_definitions/definition_topological_ring.tex}
\begin{definition}[Topological groups]  \label{definition:topological_group}
    A \hldef{topological group} is a \CrefAndHyperrefIfExist{definition:group}{group} $(G,\cdot)$ together with a \CrefAndHyperrefIfExist{definition:topological_space}{topology} $\mathcal{T}$ on $G$ such that the maps
    \begin{align*}
    \mu : G \times G &\to G, & (g,h) &\mapsto g \cdot h, \\
    \iota : G &\to G, & g &\mapsto g^{-1},
    \end{align*}
    are \CrefAndHyperrefIfExist{definition:continuous_map_between_open_subsets_of_euclidean_spaces}{continuous} with respect to the \CrefAndHyperrefIfExist{definition:product_topology_on_a_cartesian_product_of_topological_spaces}{product topology} on $G \times G$ and the topology $\mathcal{T}$ on $G$.

    A \hldef{abelian topological group} or a \hldef{commutative topological group} is a topological group that is \CrefAndHyperrefIfExist{definition:group}{abelian as a group}. \TextIfExists{definition:topological_module_over_topological_rings}{Equvialently, an abelian topological group is a \CrefAndHyperrefIfExist{definition:topological_module_over_topological_rings}{topological module} over $\bbZ$, where $\bbZ$ is given the structure of a \CrefAndHyperrefIfExist{definition:topological_ring}{topological ring} with the \CrefAndHyperrefIfExist{definition:discrete_topological_space_of_a_set}{discrete topology}.}
\end{definition}



\begin{definition} \label{definition:topological_module_over_topological_rings}
Let $R$ and $S$ be \CrefAndHyperrefIfExist{definition:topological_ring}{topological rings}. 
A \hldef{topological $(R, S)$-bimodule} is an \CrefAndHyperrefIfExist{definition:module_of_a_ring}{$(R, S)$-bimodule} $M$ equipped with a \CrefAndHyperrefIfExist{definition:topological_space}{topology} such that:
\begin{enumerate}
    \item $M$ is a \CrefAndHyperrefIfExist{definition:topological_group}{topological abelian group} under addition.
    \item The left action map $R \times M \to M$, defined by $(r, m) \mapsto rm$, is \CrefAndHyperrefIfExist{definition:continuous_map_of_topological_spaces}{continuous}.
    \item The right action map $M \times S \to M$, defined by $(m, s) \mapsto ms$, is continuous.
\end{enumerate}
Here, the product sets $R \times M$ and $M \times S$ are endowed with the \CrefAndHyperrefIfExist{definition:product_topology_on_a_cartesian_product_of_topological_spaces}{product topology}.

One also calls a topological module a \hldef{continuous module}.


A \hldef{left topological $R$-module} is a topological $(R, \mathbb{Z})$-bimodule, where $\mathbb{Z}$ carries the \CrefAndHyperrefIfExist{definition:discrete_topological_space_of_a_set}{discrete topology}.  Similarly, a \hldef{right topological $R$-module} is a topological $(\mathbb{Z}, R)$-bimodule, where $\mathbb{Z}$ carries the discrete topology. Furthemore, a \hldef{two-sided topological $R$-module} is a topological $(R, R)$-bimodule.

The discrete topology on $\mathbb{Z}$ ensures that the required continuity of the $\mathbb{Z}$-action is satisfied by any topological abelian group $M$, as the map $\mathbb{Z} \times M \to M$ is continuous if and only if for each $n \in \mathbb{Z}$, the map $m \mapsto nm$ is continuous.
\end{definition}
\begin{definition} \label{definition:morphism_of_topological_modules_over_topological_rings}
Let $R$ and $S$ be \CrefAndHyperrefIfExist{definition:topological_ring}{topological rings}, and let $M$ and $N$ be \CrefAndHyperrefIfExist{definition:topological_module_over_topological_rings}{topological $(R, S)$-bimodules}. A \hldef{morphism of topological $(R, S)$-bimodules} is a map $f: M \to N$ such that:
\begin{enumerate}
    \item $f$ is a \CrefAndHyperrefIfExist{definition:homomorphism_of_modules_over_a_ring}{morphism of $(R, S)$-bimodules} in the algebraic sense: $f(rm + m's) = r f(m) + f(m')s$ for all $r \in R, s \in S$, and $m, m' \in M$.
    \item $f$ is \CrefAndHyperrefIfExist{definition:continuous_map_of_topological_spaces}{continuous} with respect to the topologies on $M$ and $N$.

\end{enumerate}


    A \hldef{morphism of left topological $R$-modules} is a morphism of topological $(R, \mathbb{Z})$-bimodules where $\bbZ$ is endowed with the \CrefAndHyperrefIfExist{definition:discrete_topological_space_of_a_set}{discrete topology}. Similarly, a \hldef{morphism of right topological $R$-modules} is a morphism of topological $(\mathbb{Z}, R)$-bimodules where $\bbZ$ is endowed with the discrete topology. Furthermore, a \hldef{morphism of two-sided topological $R$-modules} is a morphism of topological $(R, R)$-bimodules.

    In the case of a left topological $R$-module, the condition of being a $(\mathbb{Z}, \text{discrete})$-bimodule morphism reduces to $f$ being an $R$-linear continuous map of topological abelian groups, as the $\mathbb{Z}$-linearity $f(nz) = nf(z)$ is automatically satisfied by any continuous homomorphism of topological abelian groups.
\end{definition}
\begin{definition} \label{definition:categories_of_topological_modules_over_topological_rings}
Let $R$ and $S$ be topological rings. Equip $\mathbb{Z}$ with the discrete topology.
\begin{enumerate}
    \item The \hldef{category of topological $(R, S)$-bimodules}, denoted \hl{${}_{R}\mathbf{TopMod}_{S}$}, is the category whose objects are topological $(R, S)$-bimodules and whose morphisms are continuous $(R, S)$-bimodule homomorphisms.

    \item The category of \hldef{left topological $R$-modules}, denoted \hl{$R\text{-}\mathbf{TopMod}$}, is the category ${}_{R}\mathbf{TopMod}_{\mathbb{Z}}$.

    \item The category of \hldef{right topological $R$-modules}, denoted \hl{$\mathbf{TopMod}\text{-}R$}, is the category ${}_{\mathbb{Z}}\mathbf{TopMod}_{R}$.
\end{enumerate}

\end{definition}

\begin{proposition}
An object $M \in R\text{-}\mathbf{TopMod}$ is precisely a topological abelian group equipped with a continuous left $R$-action that is distributive over addition. A morphism $f: M \to N$ in $R\text{-}\mathbf{TopMod}$ is a continuous $R$-linear map.
\end{proposition}


\begin{proposition}
Let $M$ be an $R$-module. Endowing $M$ with a topology makes it a left topological $R$-module if and only if addition is continuous and the map $R \times M \to M$ is continuous. The right action of $(\mathbb{Z}, \text{discrete})$ is then automatically continuous if and only if addition is continuous, as the action $m \cdot n$ is simply iterated addition or its inverse.
\end{proposition}


\input{../_definitions/definition_topological_module_over_a_group_ring_of_a_topological_group_over_a_topological_ring.tex}

\begin{definition} \label{definition:morphism_of_topological_modules_of_a_group_ring}
Let $R$ be a \CrefAndHyperrefIfExist{definition:topological_ring}{topological ring} and $G$ be a \CrefAndHyperrefIfExist{definition:topological_group}{topological group}. 
A \hldef{morphism of topological $G$-modules over $R$} is an \CrefAndHyperrefIfExist{definition:homomorphism_of_modules_over_a_ring}{$R$-linear} map $f: M \to N$ that is both a \CrefAndHyperrefIfExist{definition:continuous_map_of_topological_spaces}{continuous function} of topological spaces, and a $G$-equivariant map (i.e., $f(g \cdot m) = g \cdot f(m)$ for all $g \in G, m \in M$).
\end{definition}

\begin{definition} \label{definition:category_of_topological_modules_over_a_group_ring_of_a_topological_group_over_a_topological_ring}
    Let $R$ be a \CrefAndHyperrefIfExist{definition:topological_ring}{topological ring} and $G$ a \CrefAndHyperrefIfExist{definition:topological_group}{topological group}. 
    The \hldef{category of (left) topological $G$-modules over $R$} is the \CrefAndHyperrefIfExist{definition:category}{category} \hl{$R\text{-}\mathbf{TopMod}_G$} whose objects are \CrefAndHyperrefIfExist{definition:topological_module_over_a_group_ring_of_a_topological_group_over_a_topological_ring}{(left) topological $G$-modules over $R$} and whose morphisms are \CrefAndHyperrefIfExist{definition:morphism_of_topological_modules_of_a_group_ring}{morphisms of topological $G$-modules over $R$}.

    Note that we may symmetrically define the \hldef{category of right topological $G$-modules over $R$}, which may be notated by notations such as \hl{${}_G\mathbf{TopMod}-R$}.
\end{definition}

\subsection{Continuous cohomology of a topological module over a group ring of a topological group over a topological ring}


\begin{definition} \label{definition:G_invariant_functor_of_topological_group_module_over_topological_ring}
Let $R$ be a \CrefAndHyperrefIfExist{definition:topological_ring}{topological ring} and $G$ a \CrefAndHyperrefIfExist{definition:topological_group}{topological group}. Let $R\text{-}\mathbf{TopMod}_G$ be the \CrefAndHyperrefIfExist{definition:category_of_topological_modules_over_a_group_ring_of_a_topological_group_over_a_topological_ring}{category of (left) topological $G$-modules over $R$}.

The \hldef{$G$-invariant functor} \hl{$(-)^G: R\text{-}\mathbf{TopMod}_G \to R\text{-}\mathbf{TopMod}$} (\Cref{definition:categories_of_topological_modules_over_topological_rings}) is defined by:
$$M^G = \{ m \in M \mid g \cdot m = m, \, \forall g \in G \}.$$
In other words, $M^G$ is simply the submodule of \CrefAndHyperrefIfExist{definition:invariant_subgroup_of_a_group_acting_on_an_abelian_group}{invariants of $M$} as an \CrefAndHyperrefIfExist{definition:left_right_module_over_a_group_ring_of_a_group_over_a_ring}{$R[G]$-module} without considering the topological structures.
\end{definition}
\begin{definition} \label{definition:continuous_cochain_complex_of_a_topological_module_over_a_group_ring_of_a_topological_group_over_a_topological_ring}
Let $R$ be a \CrefAndHyperrefIfExist{definition:topological_ring}{topological ring} and $G$ a \CrefAndHyperrefIfExist{definition:topological_group}{topological group}. Let $R\text{-}\mathbf{TopMod}_G$ be the \CrefAndHyperrefIfExist{definition:category_of_topological_modules_over_a_group_ring_of_a_topological_group_over_a_topological_ring}{category of (left) topological $G$-modules over $R$}.

Let $M$ be an object in $R\text{-}\mathbf{TopMod}_G$. The \hldef{continuous cochain complex} \hl{$C^\bullet_{\text{cont}}(G, M)$} is defined as follows:
\begin{enumerate}
    \item $C^n_{\text{cont}}(G, M) = \text{Map}_{\text{cont}}(G^n, M)$, the $R$-module of continuous maps from $G^n$ to $M$.
    \item The \hldef{coboundary map} \hl{$d^n: C^n_{\text{cont}}(G, M) \to C^{n+1}_{\text{cont}}(G, M)$} is given by:
    $$(d^n f)(g_1, \dots, g_{n+1}) = g_1 f(g_2, \dots, g_{n+1}) + \sum_{i=1}^n (-1)^i f(g_1, \dots, g_i g_{i+1}, \dots, g_{n+1}) + (-1)^{n+1} f(g_1, \dots, g_n)$$
\end{enumerate}
\end{definition}
\begin{definition} \label{definition:continuous_cohomology_groups_of_a_topological_module_over_a_topological_group_and_a_topological_ring}

    % \TODO{It is probably more accurate to a priori define continuous cohomology by the chain complexes and then define it as the derived functor assuming that enough injectives existe}


Let $R$ be a \CrefAndHyperrefIfExist{definition:topological_ring}{topological ring} and $G$ a \CrefAndHyperrefIfExist{definition:topological_group}{topological group}. Let $R\text{-}\mathbf{TopMod}_G$ be the \CrefAndHyperrefIfExist{definition:category_of_topological_modules_over_a_group_ring_of_a_topological_group_over_a_topological_ring}{category of (left) topological $G$-modules over $R$}.

The \hldef{continuous cohomology groups} \hl{$H^n_{\text{cont}}(G, M)$}  are defined as the cohomology groups of the \CrefAndHyperref{definition:continuous_cochain_complex_of_a_topological_module_over_a_group_ring_of_a_topological_group_over_a_topological_ring}{continuous cochain complex} $C^\bullet_{\text{cont}}(G, M)$.

The continuous cohomology group is also denoted by \hl{$H^n(G, M)$} if the continuity is clear from context, but this should not be confused for the \CrefAndHyperrefIfExist{definition:group_cohomology_homology_of_a_module_over_a_group_ring_over_a_ring}{(non-continuous) cohomology group} $H^n(G,M)$ of $M$ as a $G$-module (without consideration of the topological structures).

\end{definition}
\begin{theorem} \label{theorem:continuous_cohomology_of_topological_module_over_topological_gropu_and_topological_ring_can_be_computed_by_continuous_cohomology_groups}
    \TODO{verify these facts}

Let $R$ be a \CrefAndHyperrefIfExist{definition:topological_ring}{topological ring} and $G$ a \CrefAndHyperrefIfExist{definition:topological_group}{topological group}. Let $R\text{-}\mathbf{TopMod}_G$ be the \CrefAndHyperrefIfExist{definition:category_of_topological_modules_over_a_group_ring_of_a_topological_group_over_a_topological_ring}{category of (left) topological $G$-modules over $R$}.

Consider the \CrefAndHyperrefIfExist{definition:G_invariant_functor_of_topological_group_module_over_topological_ring}{$G$-invariant functor} $(-)^G: R\text{-}\mathbf{TopMod}_G \to R\text{-}\mathbf{TopMod}$ (\Cref{definition:categories_of_topological_modules_over_topological_rings}).

Assume either of the following:
\begin{enumerate}
    \item $G$ is \CrefAndHyperrefIfExist{definition:profinite_group}{profinite}, $R$ is discrete, and $M$ is \CrefAndHyperrefIfExist{definition:discrete_group}{discrete}; in this case, by the ``right derived functors of $(-)^G$'', we mean the right derived functor on the category of discrete $G$-modules, which has enough injectives \TODO{cite; apparently the category of discrete $G$-modules over a topological ring $R$ where $G$ is profinite is an abelian category with enough injectives}.
    \item $G$ is locally compact, and $M$ is complete; in this case, by the ``right derived functors of $(-)^G$'', we mean the right derived functor in the sense of relative homological algebra, using resolutions that split as sequences of topological spaces \TODO{figureo ut what this means}.
\end{enumerate}

The \CrefAndHyperrefIfExist{definition:continuous_cohomology_groups_of_a_topological_module_over_a_topological_group_and_a_topological_ring}{continuous cohomology groups} $H^n_{\text{cont}}(G,M)$, i.e. the  \CrefAndHyperrefIfExist{definition:homology_and_cohomology_objects_for_a_chain_complex_in_an_additive_category}{cohomology} of the \CrefAndHyperrefIfExist{definition:continuous_cochain_complex_of_a_topological_module_over_a_group_ring_of_a_topological_group_over_a_topological_ring}{continuous cochain complex} $C^\bullet_{\text{cont}}(G, M)$, is naturally isomorphic to the the right derived functors of $(-)^G$. In other words,
$$H^n(C^\bullet_{\text{cont}}(G, M)) \cong (R^n(-^G))(M)$$
\end{theorem}


In particular, Galois cohomology groups are a instances of continuous cohomology groups.


\begin{definition} \label{definition:galois_cohomology_of_a_galois_field_extension_with_coefficients_in_a_topological_galois_module_over_a_topological_ring}
    Let $R$ be a \CrefAndHyperrefIfExist{definition:topological_ring}{topological ring} (e.g., a $\ell$-adic \CrefAndHyperrefIfExist{definition:ring}{ring} $\mathbb{Z}_\ell$) \TODO{ell-adic ring}.
    \begin{enumerate}
        \item Let $L/K$ be a \CrefAndHyperrefIfExist{definition:galois_extension_of_fields}{Galois field extension}.   Let $M$ be a \CrefAndHyperrefIfExist{definition:topological_module_over_a_group_ring_of_a_topological_gropu_over_a_topological_ring}{topological $G$-module} over $R$

        The \hldef{Galois cohomology group of $L/K$ with coefficients in $M$} is:
        $$\hlin{H^n(L/K, M) = H^n_{\text{cont}}(\Gal(L/K), M)}$$
        (\Cref{definition:continuous_cohomology_groups_of_a_topological_module_over_a_topological_group_and_a_topological_ring}).

        \TODO{absolute galois group}
        \item Let $K$ be a \CrefAndHyperrefIfExist{definition:field}{field} with absolute Galois group $G_K = \Gal(\overline{K}/K)$. Let $M$ be a \CrefAndHyperrefIfExist{definition:topological_module_over_a_group_ring_of_a_topological_gropu_over_a_topological_ring}{topological $G$-module} over $R$.

        The \hldef{Galois cohomology group of $K$ with coefficients in $M$} is:
        $$\hlin{H^n(K, M) = H^n(\overline{K}/K, M) = H^n_{\text{cont}}(\Gal(\overline{K}/K), M)}.$$
    \end{enumerate}

This framework is used to define $\ell$-adic cohomology, where $M$ is often a $\mathbb{Z}_\ell$-module or a $\mathbb{Q}_\ell$-vector space arising from the Tate module of an abelian variety or the étale cohomology of a variety.
\end{definition}

\begin{proposition}
    \TODO{is this true of more general profinite groups}
If $M = \varprojlim M_i$ is an inverse limit of finite discrete $G$-modules $M_i$, the Galois cohomology groups for $n=1$ (and under certain conditions for $n > 1$) satisfy:
$$H^n_{\text{cont}}(G, \varprojlim M_i) \cong \varprojlim H^n(G, M_i)$$
where the right-hand side uses the standard cohomology for discrete modules.
\end{proposition}

\begin{theorem}
    \TODO{}
Let $K$ be a field and $M$ be a topological $G_K$-module over a topological ring $R$. If $0 \to A \to B \to C \to 0$ is a sequence of topological $G_K$-modules that is topologically split as a sequence of $R$-modules, there exists a long exact sequence of $R$-modules:
$$\dots \to H^n(K, A) \to H^n(K, B) \to H^n(K, C) \to H^{n+1}(K, A) \to \dots$$
\end{theorem}


\subsection{Cohomological dimensions of groups}


\begin{definition} \label{definition:cohomological_dimension_of_a_group_over_a_ring}
Let $G$ be a \CrefAndHyperrefIfExist{definition:group}{group}. 
The \hldef{cohomological dimension of $G$ over a ring $R$}, denoted \hl{$\text{cd}_R(G)$}, is the supremum of the set of integers $n$ such that there exists an \CrefAndHyperrefIfExist{definition:left_right_module_over_a_group_ring_of_a_group_over_a_ring}{$R[G]$-module} $M$ with $H^n(G, M) \neq 0$. If no such maximum exists, we say $\text{cd}_R(G) = \infty$. When $R = \mathbb{Z}$, we simply write \hl{$\text{cd}(G)$}.
\end{definition}


\begin{definition} \label{definition:continuous_cohomological_dimension_of_a_topological_group_over_a_topological_ring}
Let $G$ be a \CrefAndHyperrefIfExist{definition:topological_group}{topological group} and $R$ a \CrefAndHyperrefIfExist{definition:topological_ring}{topological ring}. 
\TODO{technically, need to assume that the category of $G$-modules over $R$ has enough injectives to talk about continuous cohomology groups probably}
The \hldef{continuous cohomological dimension of $G$ over $R$}, denoted \hl{$\text{cd}_R(G)$}, is the supremum of the set of integers $n$ such that there exists a \CrefAndHyperrefIfExist{definition:topological_module_over_a_group_ring_of_a_topological_gropu_over_a_topological_ring}{topological $G$-module} $M$ over $R$ for which \CrefAndHyperrefIfExist{definition:continuous_cohomology_groups_of_a_topological_module_over_a_topological_group_and_a_topological_ring}{$H^n_{\text{cont}}(G, M) \neq 0$}.
\end{definition}


\begin{definition} \label{definition:p_cohomological_dimension_of_a_topological_group}
Let $G$ be a \CrefAndHyperrefIfExist{definition:topological_group}{topological group}. 

Given a prime number $p$, the \hldef{$p$-cohomological dimension of $G$}, denoted \hl{$\text{cd}_p(G)$}, is the infimum of the set of integers $n$ such that \CrefAndHyperrefIfExist{definition:continuous_cohomology_groups_of_a_topological_module_over_a_topological_group_and_a_topological_ring}{$H^i_{\text{cont}}(G, M) = 0$} for all $i > n$ and all discrete $G$-modules $M$ that are $p$-primary torsion abelian groups.

The \hldef{cohomological dimension of a topological group $G$} is the supremum \hl{$\text{cd}(G)$} over all primes $p$ of the $p$-cohomological dimensions:
$$\text{cd}(G) = \sup_p \{ \text{cd}_p(G) \}$$

Given a profinite group $G$, the \hldef{cohomological dimension of $G$} refers to $\text{cd}_p(G)$ or $\text{cd}(G)$, as opposed to the notion of continuous cohomological dimension of \Cref{definition:continuous_cohomological_dimension_of_a_topological_group_over_a_topological_ring}.
\end{definition}



\begin{remark}
In the case of a discrete group $G$, these definitions recover the classical algebraic cohomological dimension by taking all modules to have the discrete topology. For locally compact groups, such as Lie groups, one often restricts the coefficients $M$ to a specific class, such as Fréchet spaces or Banach spaces, leading to the notion of continuous or bounded cohomological dimension.
\end{remark}


\begin{proposition}
Let $G$ be a topological group. If $G$ contains a closed subgroup $H$, then $\text{cd}_p(H) \leq \text{cd}_p(G)$. If $G$ is a profinite group and $U$ is an open subgroup, then $\text{cd}_p(U) = \text{cd}_p(G)$, a result that does not hold for arbitrary topological groups without further compactness or finiteness assumptions.
\end{proposition}


% \begin{definition}
% Let $p$ be a prime number and $G$ a profinite group. 
% The \hldef{$p$-cohomological dimension of $G$}, denoted \hl{$\text{cd}_p(G)$}, is the infimum of the set of integers $n$ such that for every discrete periodic $G$-module $M$ whose elements have order a power of $p$, we have $H^i_{\text{cont}}(G, M) = 0$ for all $i > n$.
% \end{definition}

% \begin{definition}
% The \hldef{cohomological dimension of a profinite group $G$} is defined as:
% $$\text{cd}(G) = \sup_p \{ \text{cd}_p(G) \}$$
% where the supremum is taken over all prime numbers $p$.
% \end{definition}


\begin{proposition}
The cohomological dimension of a group $G$ over $R$ is equal to the projective dimension of $R$ as a left $R[G]$-module, where $R$ is equipped with the trivial $G$-action:
$$\text{cd}_R(G) = \text{pd}_{R[G]}(R)$$
\textit{Remark:} \TODO{define projective dimension}
\end{proposition}

\begin{theorem}
Let $G$ be a profinite group and $H$ a closed subgroup. Then:
\begin{enumerate}
    \item $\text{cd}_p(H) \leq \text{cd}_p(G)$.
    \item If $H$ is an open subgroup, or more generally if $H$ is the kernel of a continuous homomorphism to a finite group of order prime to $p$, then $\text{cd}_p(H) = \text{cd}_p(G)$.
\end{enumerate}
\end{theorem}

\subsubsection{Cohomological dimensions of fields}

\begin{definition}  \label{definition:p_cohomological_dimension_of_a_field}
    Let $K$ be a \CrefAndHyperrefIfExist{definition:field}{field}. 

    Given a prime number $p$, the $p$-cohomological dimension of $K$, \hl{$\text{cd}_p(K)$}, is defined as the \CrefAndHyperref{definition:p_cohomological_dimension_of_a_topological_group}{$p$-cohomological dimension} of its absolute Galois group $G_K$.
    \TODO{absolute galois group}
    
    Similarly, the \hldef{cohomological dimension of $K$}, denoted \hl{$\text{cd}(K)$}, is defined as the \CrefAndHyperref{definition:p_cohomological_dimension_of_a_topological_group}{cohomological dimension} of $G_K$.
\end{definition}
\begin{theorem} \label{theorem:cohomological_dimensions_of_common_fields}
Let $K$ be a \CrefAndHyperrefIfExist{definition:field}{field} and $\text{cd}(K)$ its \CrefAndHyperrefIfExist{definition:p_cohomological_dimension_of_a_field}{cohomological dimension}.
\begin{enumerate}
    \item If $K$ is \CrefAndHyperrefIfExist{definition:algebraically_closed_field}{algebraically closed} or, more generally, $K$ is a \CrefAndHyperrefIfExist{definition:separably_closed_field}{separably closed field}, then $\text{cd}(K) = 0$.
    \item If $K$ is a \CrefAndHyperrefIfExist{definition:finite_field}{finite field} or a \TODO{maximal extension of a finite field} (such as the field of all roots of unity over a finite field), then $\text{cd}(K) = 1$.

    \item If $K$ is a \CrefAndHyperrefIfExist{definition:local_field}{local field} of \CrefAndHyperrefIfExist{definition:characteristic_of_a_ring}{characteristic zero} (a finite extension of $\mathbb{Q}_p$), then $\text{cd}(K) = 2$.

    \TODO{totally imaginary number field}
    \item If $K$ is a \CrefAndHyperrefIfExist{global_field}{number field} and $K$ has no real embeddings (is totally imaginary), then $\text{cd}(K) = 2$. 

    \TODO{absolute galois group}
    \item If $K$ is a number field with at least one real embedding, then $\text{cd}(K) = \infty$. This is due to the presence of real closed subfields which contain an element of order 2 in their absolute Galois group, implying $\text{cd}_2(K) = \infty$.
\end{enumerate}
\end{theorem}
\begin{theorem}[Tsen's Theorem] \label{theorem:tsens_theorem_on_the_cohomological_dimension_of_a_function_field_over_a_field}
Let $K$ be a \CrefAndHyperrefIfExist{definition:field}{field} of \CrefAndHyperrefIfExist{definition:p_cohomological_dimension_of_a_field}{cohomological dimension} $\text{cd}(K) = d$.

If $L$ is a \CrefAndHyperrefIfExist{definition:function_field_over_a_field}{function field} in $n$ variables over $K$, then:
$$\text{cd}(L) \leq d + n$$
In particular, if $K$ is an \CrefAndHyperrefIfExist{definition:algebraically_closed_field}{algebraically closed field} and $L$ is the function field of an \CrefAndHyperrefIfExist{definition:algebraic_curve_over_a_field}{algebraic curve} over $K$, then $\text{cd}(L) = 1$.
\end{theorem}
\begin{proposition} \label{proposition:p_cohomological_dimension_of_positive_characteristic_field_is_at_most_one}
Let $K$ be a \CrefAndHyperrefIfExist{definition:field}{field} of characteristic $p > 0$. 
Then $\text{cd}_p(K) \leq 1$ (\Cref{definition:p_cohomological_dimension_of_a_field}). This bound is independent of the \CrefAndHyperrefIfExist{definition:transcendence_degree_of_a_field_extension}{transcendence degree} of the field over its \CrefAndHyperrefIfExist{definition:prime_subfield_of_a_field}{prime field}.
\end{proposition}
\begin{theorem} \label{theorem:field_of_formal_laurent_series_over_algebraically_closed_field_has_cohomological_dimension_equal_to_number_of_variables}
    \TODO{laurent series}
Let $k$ be an \CrefAndHyperrefIfExist{definition:algebraically_closed_field}{algebraically closed field}. The \CrefAndHyperrefIfExist{definition:p_cohomological_dimension_of_a_field}{cohomological dimension} of the field of formal Laurent series $K = k((t_1))\dots((t_n))$ is $n$.
\end{theorem}



% \begin{definition}
% A \hldef{profinite group} $G$ is a topological group that is the inverse limit of a system of finite discrete groups.

% A \hldef{discrete $G$-module} is an abelian group $A$ with the discrete topology such that the action map $G \times A \to A$ is continuous. The \hldef{continuous cohomology} \hl{$H^n_{cont}(G, A)$} is the $n$-th right derived functor of the fixed-point functor in the category of discrete $G$-modules.

% \end{definition}

\begin{definition}
Let $\mathfrak{g}$ be a Lie algebra over a field $k$ and $M$ a $\mathfrak{g}$-module. The \hldef{Chevalley-Eilenberg complex} \hl{$C^n(\mathfrak{g}, M) = \text{Hom}_k(\Lambda^n \mathfrak{g}, M)$} consists of alternating $n$-linear maps. The \hldef{Lie algebra cohomology} \hl{$H^n(\mathfrak{g}, M)$} is the cohomology of this complex with differential \hl{$d: C^n(\mathfrak{g}, M) \to C^{n+1}(\mathfrak{g}, M)$} given by:
$$(df)(x_1, \dots, x_{n+1}) = \sum_{i < j} (-1)^{i+j} f([x_i, x_j], x_1, \dots, \hat{x}_i, \dots, \hat{x}_j, \dots, x_{n+1}) + \sum_i (-1)^{i+1} x_i \cdot f(x_1, \dots, \hat{x}_i, \dots, x_{n+1})$$
\end{definition}

\begin{definition}
Let $G$ be a Lie group and $A$ a topological $G$-module. The \hldef{differentiable cohomology} \hl{$H^n_{diff}(G, A)$} is defined as the cohomology of the complex of smooth cochains $C^\infty(G^n, A)$. \textit{Remark: \TODO{define van Est isomorphism}}.
\end{definition}

\begin{definition}
Let $S$ be a $k$-algebra and $M$ an $S$-bimodule. The \hldef{Hochschild cohomology} \hl{$HH^n(S, M)$} is defined as $\text{Ext}^n_{S \otimes_k S^{op}}(S, M)$. In the case of a group ring $S = k[G]$, this corresponds to the cohomology of $G$ acting on $M$ by conjugation.
\end{definition}

\begin{definition}
Let $G$ be a discrete group and $A$ a \hldef{Banach $G$-module} (a Banach space with an isometric $G$-action). The \hldef{bounded cohomology} \hl{$H^n_b(G, A)$} is the cohomology of the complex of bounded cochains $C^n_b(G, A) = \{f: G^n \to A \mid \sup \|f(g_1, \dots, g_n)\| < \infty\}$ with the standard bar differential.
\end{definition}

\begin{definition}
Let $G$ be an affine group scheme over a field $k$ and $M$ a rational $G$-module. \textit{Remark: \TODO{define affine group scheme}}. The cohomology \hl{$H^n(G, M)$} is defined as the Hochschild cohomology of the coordinate ring $k[G]$ in the category of $k[G]$-comodules.
\end{definition}
